% !TEX program = pdflatex
\documentclass[11pt]{article}

\usepackage[T1]{fontenc}
\usepackage{lmodern}
\usepackage{microtype}
\usepackage[a4paper,margin=28mm,headheight=14pt]{geometry}
\usepackage{amsmath,amssymb,amsthm,mathtools}
\usepackage{enumitem}
\usepackage{array,booktabs,longtable}
\usepackage{xcolor}
\usepackage{xurl}
\usepackage{hyperref}
\usepackage{fancyhdr}

\definecolor{linkblue}{RGB}{24,78,119}
\definecolor{citegreen}{RGB}{36,104,71}
\hypersetup{
  colorlinks=true,
  linkcolor=linkblue,
  citecolor=citegreen,
  urlcolor=linkblue,
  pdftitle={A Dirichlet Polya Inequality for Three-Dimensional Spherical Shells},
  pdfsubject={Purely analytic main proof manuscript},
  pdfkeywords={Polya conjecture, Dirichlet Laplacian, spherical shell, eigenvalue counting}
}

\pagestyle{fancy}
\fancyhf{}
\fancyhead[L]{\small Dirichlet P\'olya inequality for spherical shells}
\fancyhead[R]{\small Purely analytic proof}
\fancyfoot[C]{\thepage}
\fancypagestyle{plain}{%
  \fancyhf{}%
  \fancyhead[L]{\small Dirichlet P\'olya inequality for spherical shells}%
  \fancyhead[R]{\small Purely analytic proof}%
  \fancyfoot[C]{\thepage}%
}
\fancypagestyle{titlepage}{%
  \fancyhf{}%
  \fancyfoot[C]{\thepage}%
  \renewcommand{\headrulewidth}{0pt}%
}

\numberwithin{equation}{section}
\allowdisplaybreaks[2]
\setlist{itemsep=0.25em,topsep=0.45em}
\setlength{\emergencystretch}{3em}
\newtheorem{theorem}{Theorem}[section]
\newtheorem{proposition}[theorem]{Proposition}
\newtheorem{lemma}[theorem]{Lemma}
\newtheorem{corollary}[theorem]{Corollary}
\theoremstyle{definition}
\newtheorem{definition}[theorem]{Definition}
\theoremstyle{remark}
\newtheorem{remark}[theorem]{Remark}

\newcommand{\R}{\mathbb R}
\newcommand{\N}{\mathbb N}
\newcommand{\Nzero}{\mathbb N_0}
\DeclareMathOperator{\arsinh}{arsinh}

\title{A Dirichlet P\'olya Inequality for\\Three-Dimensional Spherical Shells\\[0.4em]
\large Purely analytic main proof manuscript}
\author{Anonymous for peer review}
\date{17 July 2026}

\begin{document}

\maketitle
\thispagestyle{titlepage}

\begin{quote}
\small
\textbf{Scope and status.}
This manuscript proves a theorem for the class of genuine three-dimensional
spherical shells.  It does not claim the general P\'olya conjecture,
literature novelty, or priority.  All shell-specific analytic reductions,
regional arguments, boundary assignments, exact constants, and all finite
exact-rational decisions are stated in this paper and its accompanying
purely analytic supplement.  Executable replay sources and interval
certificates are retained only as optional regression checks and are not
premises of the proof.  The imported background and analytic results
comprise explicitly identified standard Sturm--Liouville and Bessel facts and
published phase, Bessel-zero, and one-dimensional floor-sum inequalities,
each quoted with its precise scope.
\end{quote}

\begin{abstract}
For \(0<r<R\), let
\[
  A_{r,R}=\{x\in\R^3:r<|x|<R\}.
\]
We prove, with the strict Dirichlet counting convention, that
\[
  N_D(A_{r,R},\Lambda)
  \le L_3|A_{r,R}|\Lambda^{3/2},
  \qquad
  L_3=\frac1{6\pi^2},
  \qquad \Lambda\ge0.
\]
The proof is purely analytic.  It separates variables, converts the radial
spectrum into a strict multiplicity-weighted phase count, and combines an
exact retained-remainder identity with a ratio-sharp angular block estimate,
a tangent-envelope radial minorant, a universal ball quarter-layer, and an
elementary scalar convexity argument.  The low- and middle-ratio proof starts
at the first possible radial frequency; no finite Bessel-zero staircase,
event-state theorem, floating-point conclusion, or interval-arithmetic
conclusion is a proof premise.  We also give an elementary proof that no genuine spherical
shell tiles \(\R^3\) by congruent rigid-motion copies with disjoint interiors,
under either exact or almost-everywhere coverage.
\end{abstract}

\tableofcontents
\clearpage

\section{Main results and conventions}

The eigenvalues of the Dirichlet Laplacian on a bounded domain \(\Omega\) are
listed with multiplicity as
\[
  0<\lambda_1^D(\Omega)\le\lambda_2^D(\Omega)\le\cdots.
\]
Throughout we use the strict counting function
\begin{equation}
  N_D(\Omega,\Lambda)
  :=\#\{j:\lambda_j^D(\Omega)<\Lambda\}.
  \label{eq:strict-count}
\end{equation}
This convention matters at every spectral wall.  We write
\(\N=\{1,2,3,\ldots\}\) and \(\Nzero=\{0,1,2,\ldots\}\).

\begin{theorem}[Spectral inequality]
\label{thm:spectral}
For every \(0<r<R\) and every \(\Lambda\ge0\),
\begin{equation}
  \boxed{
  N_D(A_{r,R},\Lambda)
  \le \frac{2}{9\pi}(R^3-r^3)\Lambda^{3/2}
  =L_3|A_{r,R}|\Lambda^{3/2},
  \qquad L_3=\frac1{6\pi^2}.}
  \label{eq:main-theorem}
\end{equation}
\end{theorem}

The standard non-strict formulation follows without changing the constant.
Indeed, if
\[
 N_D^{\leq}(\Omega,\Lambda)
 :=\#\{j:\lambda_j^D(\Omega)\leq\Lambda\},
\]
then discreteness of the Dirichlet spectrum gives
\begin{equation}
 N_D^{\leq}(\Omega,\Lambda)
 =\lim_{\varepsilon\downarrow0}
 N_D(\Omega,\Lambda+\varepsilon).
 \label{eq:strict-to-nonstrict}
\end{equation}
Applying Theorem~\ref{thm:spectral} at
\(\Lambda+\varepsilon\) and passing to the limit proves
\eqref{eq:main-theorem} with \(N_D^{\leq}\) as well.  Thus the strict
convention used for the wall bookkeeping is equivalent here to the usual
non-strict P\'olya inequality.

\begin{theorem}[Non-tiling]
\label{thm:nontiling}
For every \(0<r<R\), no family of congruent rigid-motion copies of
\(A_{r,R}\) has pairwise-disjoint interiors and covers \(\R^3\) exactly or
up to a three-dimensional Lebesgue-null set.  The same conclusion holds when
coverage is formulated using closed shells.
\end{theorem}

The two theorems concern literally the same complete shell family.
Theorem~\ref{thm:nontiling} is logically independent of the spectral proof.

\subsection*{Proof map}

For \(K>0\), the spectral argument has one global comparison, a no-mode
region, and two high-frequency ratio owners.  After scaling to
\(A_{\rho,1}\), separation of variables and the uniform Bessel-phase
enclosure give a strict multiplicity-weighted lattice proxy \(P\) with
\[
 N_D(A_{\rho,1},K^2)\leq P.
\]
In the original action variables of Part~VII of the supplement, let \(R\) be
the inverse action, \(F=R^2\), and let \(M(r)\) count the half-integers
strictly below \(r\).  The inverse action gives the exact layer-cake picture
\[
 P=\sum_{\substack{n\geq1\\n-1/4<T}}M\bigl(R(n-1/4)\bigr)^2,
 \qquad
 W=\int_0^T R(t)^2\,dt,
\]
whose scaled version is Lemma~\ref{thin:lem-layer-cake}.  Thus the master
comparison is
\begin{equation}
 W-P=\mathcal D_{\rm rad}-\mathcal E_{\rm ang}:
 \quad\hbox{radial quadrature reserve pays angular rounding.}
 \label{eq:proof-map-master-defect}
\end{equation}
The exact definitions and the retained positive remainders in
\eqref{eq:proof-map-master-defect} are given in the accompanying
\emph{Purely Analytic Supplement}.

\begin{samepage}
The parameter space is closed by the following disjoint owners.
\begin{center}
\begin{tabular}{@{}lll@{}}
\toprule
frequency range&ratio range&owner\\
\midrule
\(0\leq K\leq\pi/(1-\rho)\)&\(0<\rho<1\)&strict no-mode bound\\
\(K>\pi/(1-\rho)\)&\(0<\rho<39/50\)&global defect theorem\\
\(K>\pi/(1-\rho)\)&\(39/50\leq\rho<1\)&optical theorem\\
\bottomrule
\end{tabular}
\end{center}
\end{samepage}
Their union proves the unit-shell inequality; dilation proves the result for
every \(A_{r,R}\).  The former finite staircase, 38-state theorem, and
611-row allocation are retained in the repository only as historical audit
material; they have no implication into the proof below.

% Self-contained manuscript fragment: separated spectrum, phase reduction,
% weighted lattice reduction, and the small-hole endpoint.
% Assumes amsmath, amssymb, amsthm, mathtools and theorem environments
% theorem, proposition, lemma.

\section{Separated spectrum and phase reduction}
\label{sec:sc-spectrum-small-hole}

Throughout this section
\[
  A_{\rho,1}=\{x\in\mathbb R^3:\rho<|x|<1\},
  \qquad 0<\rho<1,
\]
and the Dirichlet counting function is strict:
\[
  N_D(A_{\rho,1},K^2)
  :=\#\{j:\lambda_j^D(A_{\rho,1})<K^2\}.
\]
For \(x>0\) we use the strict positive-integer count
\begin{equation}
  [x]_<:=\#\{n\in\mathbb N:n<x\}
  =\max\{0,\lceil x\rceil-1\}.
  \label{eq:sc-strict-bracket}
\end{equation}
Thus \([m]_< =m-1\) for \(m\in\mathbb N\).  This convention is retained
at every eigenfrequency, phase level, and half-integer angular wall.

\subsection{The exact separated spectrum}

Let \(H_\rho=-\Delta^D_{A_{\rho,1}}\) be the Friedrichs operator associated
with
\[
  \mathfrak q[f]=\int_{A_{\rho,1}}|\nabla f|^2\,dx,
  \qquad D(\mathfrak q)=H_0^1(A_{\rho,1}).
\]
Choose an orthonormal spherical-harmonic basis satisfying
\[
  -\Delta_{\mathbb S^2}Y_{\ell m}
  =\ell(\ell+1)Y_{\ell m},
  \qquad \ell\in\mathbb N_0,
  \quad -\ell\le m\le\ell.
\]

\begin{proposition}[Unitary decomposition and completeness]
\label{prop:sc-separated-spectrum}
For
\[
  L_\ell=-\frac{d^2}{dr^2}+\frac{\ell(\ell+1)}{r^2},
  \qquad
  D(L_\ell)=H^2(\rho,1)\cap H_0^1(\rho,1),
\]
one has the unitary equivalence
\begin{equation}
  H_\rho\simeq
  \bigoplus_{\ell=0}^{\infty}
  \bigoplus_{m=-\ell}^{\ell}L_\ell.
  \label{eq:sc-operator-sum}
\end{equation}
The exact form domain on the right is
\begin{equation}
  \left\{(u_{\ell m}):
  \begin{array}{l}
  u_{\ell m}\in H_0^1(\rho,1),\\
  \displaystyle
  \sum_{\ell,m}\left(\|u_{\ell m}\|_2^2
  +\mathfrak a_\ell[u_{\ell m}]\right)<\infty
  \end{array}\right\},
  \label{eq:sc-global-form-domain}
\end{equation}
where
\[
  \mathfrak a_\ell[u]
  =\int_\rho^1\left(|u'|^2
  +\frac{\ell(\ell+1)}{r^2}|u|^2\right)dr.
\]
The exact operator domain is
\begin{equation}
  \left\{(u_{\ell m}):
  \begin{array}{l}
  u_{\ell m}\in D(L_\ell),\\
  \displaystyle
  \sum_{\ell,m}\left(\|u_{\ell m}\|_2^2
  +\|L_\ell u_{\ell m}\|_2^2\right)<\infty
  \end{array}\right\}.
  \label{eq:sc-global-operator-domain}
\end{equation}
Each \(L_\ell\) has positive, simple eigenvalues tending to infinity and a
complete orthonormal eigenbasis.  The full operator has compact resolvent.
\end{proposition}

\begin{proof}
Write \(x=r\omega\), so \(dx=r^2\,dr\,d\omega\), and set
\[
  R_{\ell m}(r)=\int_{\mathbb S^2}
  f(r,\omega)\overline{Y_{\ell m}(\omega)}\,d\omega,
  \qquad u_{\ell m}(r)=rR_{\ell m}(r).
\]
Fubini, Tonelli, and Parseval on \(\mathbb S^2\) give
\[
  \|f\|_{L^2(A_{\rho,1})}^2
  =\sum_{\ell,m}\int_\rho^1|R_{\ell m}(r)|^2r^2\,dr
  =\sum_{\ell,m}\int_\rho^1|u_{\ell m}(r)|^2\,dr.
\]
Conversely, if \((u_{\ell m})\) is square summable, the finite sums
\[
  f_L(r,\omega)=\frac1r
  \sum_{\ell\le L}\sum_{m=-\ell}^{\ell}
  u_{\ell m}(r)Y_{\ell m}(\omega)
\]
are Cauchy in \(L^2(A_{\rho,1})\), and their limit maps back to the given
family.  Hence the coefficient map is unitary onto the Hilbert direct sum.
Because \(r,1/r\in W^{1,\infty}(\rho,1)\), multiplication by either function
is bounded on \(H^1(\rho,1)\) and preserves the two zero endpoint traces.

For a finite smooth harmonic expansion, polar coordinates and angular
orthogonality yield
\[
  \mathfrak q[f]=\sum_{\ell,m}\int_\rho^1
  \left(r^2|R_{\ell m}'|^2
  +\ell(\ell+1)|R_{\ell m}|^2\right)dr.
\]
If \(u=rR\), then
\[
  r^2|R'|^2
  =\left|u'-\frac ur\right|^2
  =|u'|^2-\frac{(|u|^2)'}r+\frac{|u|^2}{r^2}.
\]
Since
\[
  \frac{(|u|^2)'}r
  =\left(\frac{|u|^2}{r}\right)'
  +\frac{|u|^2}{r^2}
\]
and \(u(\rho)=u(1)=0\), integration gives
\[
  \int_\rho^1r^2|R'|^2\,dr
  =\int_\rho^1|u'|^2\,dr.
\]
The angular term becomes
\(\ell(\ell+1)|u|^2/r^2\).  Closing finite harmonic expansions in the
form norm proves \eqref{eq:sc-global-form-domain}: one inclusion follows
from componentwise form convergence and Fatou's lemma, while the reverse
inclusion follows by first truncating the square-summable family in
\((\ell,m)\) and then approximating its finitely many components by
\(C_c^\infty(\rho,1)\).  The uniqueness theorem for operators associated
with closed semibounded forms gives \eqref{eq:sc-operator-sum}.

For fixed \(\ell\), the potential \(\ell(\ell+1)/r^2\) is continuous and
bounded on \([\rho,1]\).  The regular one-dimensional Dirichlet form
therefore has associated operator \(L_\ell\) with the displayed
\(H^2\cap H_0^1\) domain.  The standard domain formula for an orthogonal
operator sum gives \eqref{eq:sc-global-operator-domain}.

The embedding \(H_0^1(\rho,1)\hookrightarrow L^2(\rho,1)\) is compact, so
every block has compact resolvent.  To pass to the infinite direct sum, note
that \(r\le1\) implies
\[
  L_\ell\ge\ell(\ell+1),
  \qquad
  \|(L_\ell+1)^{-1}\|
  \le\frac1{1+\ell(\ell+1)}.
\]
The resolvent of the full operator is consequently the operator-norm limit
of its finite-\(\ell\) truncations, each of which is compact.  Hence the full
resolvent is compact.

Regular Dirichlet Sturm--Liouville theory gives a complete sequence of
simple eigenvalues in every block.  They are strictly positive because
\[
  \mathfrak a_\ell[u]\ge\int_\rho^1|u'|^2\,dr
  \ge\frac{\pi^2}{(1-\rho)^2}\|u\|_2^2.
\]
Simplicity also follows directly: two solutions of the same second-order
equation that vanish at \(r=\rho\) have proportional initial derivatives
and hence are proportional by uniqueness.  If \(u_{\ell,n}\) is a complete
radial eigenbasis, then
\[
  \frac{u_{\ell,n}(r)}rY_{\ell m}(\omega)
\]
over all \((\ell,m,n)\) is exactly the inverse image of a complete
orthonormal coordinate basis of the Hilbert direct sum.  This proves global
completeness.
\end{proof}

\begin{proposition}[Bessel determinant and strict phase enumeration]
\label{prop:sc-bessel-phase-count}
Put \(\nu=\ell+\tfrac12\).  The positive eigenfrequencies in the
\(\ell\)-th radial block are precisely the positive zeros of
\begin{equation}
  D_{\nu,\rho}(k)
  :=J_\nu(\rho k)Y_\nu(k)-Y_\nu(\rho k)J_\nu(k).
  \label{eq:sc-bessel-determinant}
\end{equation}
Every such zero is simple.  With the continuous phase branch
\begin{equation}
  J_\nu(x)+iY_\nu(x)=M_\nu(x)e^{i\theta_\nu(x)},
  \qquad M_\nu(x)>0,
  \qquad \theta_\nu(0+)=-\frac\pi2,
  \label{eq:sc-phase-branch}
\end{equation}
and
\[
  \Psi_{\nu,\rho}(K)=\theta_\nu(K)-\theta_\nu(\rho K),
  \qquad
  \gamma_{\rho,K}(\nu)=\frac{\Psi_{\nu,\rho}(K)}\pi,
\]
the exact strict count is
\begin{equation}
  N_D(A_{\rho,1},K^2)
  =\sum_{\ell=0}^{\infty}(2\ell+1)
  [\gamma_{\rho,K}(\ell+\tfrac12)]_<.
  \label{eq:sc-exact-phase-count}
\end{equation}
The sum is finite.  If \(K\) is a radial eigenfrequency in one or several
angular channels, the root at \(K\) is excluded in each such channel and the
excluded angular multiplicities add.
\end{proposition}

\begin{proof}
For \(k>0\), the physical radial equation has the general solution
\[
  R(r)=r^{-1/2}
  \bigl(c_1J_\nu(kr)+c_2Y_\nu(kr)\bigr).
\]
The two Dirichlet conditions admit a nonzero coefficient vector exactly when
\[
  \det\begin{pmatrix}
  J_\nu(\rho k)&Y_\nu(\rho k)\\
  J_\nu(k)&Y_\nu(k)
  \end{pmatrix}=D_{\nu,\rho}(k)=0.
\]
The Wronskian
\begin{equation}
  J_\nu(x)Y_\nu'(x)-J_\nu'(x)Y_\nu(x)=\frac2{\pi x}
  \label{eq:sc-bessel-wronskian}
\end{equation}
shows that the two entries in either row cannot vanish simultaneously.
Thus the determinant criterion remains exact even if the \(J_\nu\) values,
or the \(Y_\nu\) values, happen to vanish at both endpoints.

For root simplicity define the left-normalized solution
\begin{equation}
  s_\ell(r,k)=\frac{\pi\sqrt\rho}{2}\sqrt r
  \left[J_\nu(\rho k)Y_\nu(kr)
  -Y_\nu(\rho k)J_\nu(kr)\right].
  \label{eq:sc-left-normalized}
\end{equation}
Equation \eqref{eq:sc-bessel-wronskian} gives
\[
  s_\ell(\rho,k)=0,
  \qquad \partial_rs_\ell(\rho,k)=1,
  \qquad
  s_\ell(1,k)=\frac{\pi\sqrt\rho}{2}D_{\nu,\rho}(k).
\]
Let \(v=\partial_ks_\ell\).  Differentiating
\((L_\ell-k^2)s_\ell=0\) gives
\[
  (L_\ell-k^2)v=2ks_\ell,
  \qquad v(\rho,k)=\partial_rv(\rho,k)=0.
\]
At a positive determinant root, Lagrange's identity yields
\begin{equation}
  \partial_ks_\ell(1,k)\,\partial_rs_\ell(1,k)
  =2k\int_\rho^1s_\ell(r,k)^2\,dr>0.
  \label{eq:sc-root-simplicity}
\end{equation}
Thus \(\partial_ks_\ell(1,k)\ne0\), and the zero of the determinant is
simple.

At zero energy the left-normalized solution is
\[
  s_\ell(r,0)=
  \frac{\rho^{-\ell}r^{\ell+1}
  -\rho^{\ell+1}r^{-\ell}}{2\ell+1},
\]
so
\[
  s_\ell(1,0)=
  \frac{\rho^{-\ell}-\rho^{\ell+1}}{2\ell+1}>0.
\]
Equivalently, the standard small-argument expansions give
\begin{equation}
  \lim_{k\downarrow0}D_{\nu,\rho}(k)
  =\frac{\rho^{-\nu}-\rho^\nu}{\pi\nu}>0.
  \label{eq:sc-zero-determinant-limit}
\end{equation}
Hence the limiting phase level zero is not an eigenfrequency.

From \eqref{eq:sc-phase-branch},
\[
  J_\nu=M_\nu\cos\theta_\nu,
  \qquad Y_\nu=M_\nu\sin\theta_\nu,
\]
and direct substitution gives the sign-sensitive factorization
\begin{equation}
  D_{\nu,\rho}(k)
  =M_\nu(\rho k)M_\nu(k)
  \sin\Psi_{\nu,\rho}(k).
  \label{eq:sc-determinant-factorization}
\end{equation}
The Wronskian also gives
\[
  \theta_\nu'(x)=\frac2{\pi xM_\nu(x)^2}>0.
\]
Nicholson's integral representation \cite[Eq.~(10.9.30)]{DLMF}
\[
  M_\nu(x)^2=\frac8{\pi^2}\int_0^\infty
  \cosh(2\nu t)K_0(2x\sinh t)\,dt
\]
and the strict decrease of \(K_0\) imply that \(M_\nu(x)^2\) is strictly
decreasing in \(x\).  Therefore
\begin{equation}
  \Psi_{\nu,\rho}'(k)
  =\frac2{\pi k}
  \left(M_\nu(k)^{-2}-M_\nu(\rho k)^{-2}\right)>0.
  \label{eq:sc-phase-monotonicity}
\end{equation}
The branch condition gives \(\Psi_{\nu,\rho}(0+)=0\), while the standard
large-argument Bessel expansion gives
\[
  \Psi_{\nu,\rho}(k)=(1-\rho)k+O(k^{-1})\longrightarrow\infty
\]
for fixed \(\nu\) and \(\rho\).  Hence the positive determinant roots are
uniquely indexed by
\[
  \Psi_{\ell+1/2,\rho}(k_{\ell,n})=n\pi,
  \qquad n=1,2,\ldots.
\]
Strict monotonicity gives
\[
  k_{\ell,n}<K
  \quad\Longleftrightarrow\quad
  n\pi<\Psi_{\ell+1/2,\rho}(K).
\]
The direct-sum decomposition supplies angular multiplicity \(2\ell+1\).
If the same numerical eigenvalue occurs at distinct values of \(\ell\), the
corresponding spherical-harmonic sectors are orthogonal and their
multiplicities add.  Finally,
\(\lambda_{\ell,n}\ge\ell(\ell+1)\), so only finitely many angular channels
can contribute below \(K^2\).  These observations prove
\eqref{eq:sc-exact-phase-count}, including every strict spectral wall.
\end{proof}

The only standard inputs in the preceding proof are completeness of the
spherical harmonics, regular one-dimensional Dirichlet Sturm--Liouville
theory, the Bessel equation, the Wronskian, Nicholson's representation, and
the usual small- and large-argument Bessel formulas; see \cite{DLMF}.

\subsection{Uniform phase enclosures and the strict lattice proxy}

For \(a>0\), define
\begin{equation}
  G_a(z)=
  \begin{cases}
  \displaystyle\frac1\pi\left(\sqrt{a^2-z^2}
  -z\arccos\frac za\right),&0\le z\le a,\\[1mm]
  0,&z>a.
  \end{cases}
  \label{eq:sc-model-action}
\end{equation}

\begin{lemma}[Elementary turning bound]
\label{lem:sc-turning-bound}
For \(a>0\) and \(0\leq z<a\),
\begin{equation}
  G_a(z)<\frac{(a-z)^{3/2}}{3\sqrt a}.
  \label{eq:sc-turning-point-bound}
\end{equation}
\end{lemma}

\begin{proof}
Write \(t=\cos\theta\).  Concavity of \(\sin\) on
\([0,\pi/4]\) gives
\(\sin(\theta/2)\geq\sqrt2\,\theta/\pi\), and hence
\[
  \arccos t\leq\frac\pi2\sqrt{1-t},
  \qquad 0\leq t\leq1.
\]
Since \(G_a(a)=0\) and
\[
  G_a(z)=\frac1\pi\int_z^a\arccos\frac ta\,dt,
\]
integration gives the stated strict inequality.
\end{proof}

For \(0\le z<a\), put
\[
  H_a(z)=\frac{3a^2+2z^2}{24\pi(a^2-z^2)^{3/2}},
\]
and define
\begin{equation}
  \mathcal F_a(z)=
  \begin{cases}
  \max\{G_a(z)-H_a(z),-1/4\},&0\le z<a,\\
  -1/4,&z\ge a.
  \end{cases}
  \label{eq:sc-auxiliary-F}
\end{equation}

We use the following published Bessel-phase result, stated here with every
hypothesis needed in the sequel.

\begin{lemma}[External uniform Bessel-phase enclosure]
\label{lem:sc-external-phase}
With the branch \eqref{eq:sc-phase-branch}, for every \(a>0\) and real
\(z\ge0\),
\begin{equation}
  \mathcal F_a(z)-\frac14
  <\frac{\theta_z(a)}\pi
  <G_a(z)-\frac14.
  \label{eq:sc-individual-phase-enclosure}
\end{equation}
If \(0<\mu<K\), \(0\le z\le K\), and
\[
  \Gamma_{K,\mu}(z)
  :=\frac{\theta_z(K)-\theta_z(\mu)}\pi,
\]
then
\begin{equation}
  0<\Gamma_{K,\mu}(z)
  <G_K(z)-\mathcal F_\mu(z)
  \le G_K(z)+\frac14.
  \label{eq:sc-global-difference-enclosure}
\end{equation}
If \(0\le z\le\mu\), including \(z=\mu\), then also
\begin{equation}
  G_K(z)-G_\mu(z)
  <\Gamma_{K,\mu}(z)
  <G_K(z)-G_\mu(z)+\frac14.
  \label{eq:sc-action-difference-enclosure}
\end{equation}
No value of the divergent expression \(H_\mu(\mu)\) is used in the last
statement.
\end{lemma}

Lemma~\ref{lem:sc-external-phase} is Theorem~5.1 and Lemma~5.2 of
\cite{FLPSannuli}, based on the real-order phase enclosures of
\cite{FLPSphase}.  The upper bound in
\eqref{eq:sc-global-difference-enclosure} follows immediately by subtracting
the lower bound for \(\theta_z(\mu)\) from the upper bound for
\(\theta_z(K)\) in \eqref{eq:sc-individual-phase-enclosure}; the correlated
two-sided estimate \eqref{eq:sc-action-difference-enclosure} is the sharper
part of the annulus result.  Both theorems are stated for real order, so the
half-integers \(z=\ell+\tfrac12\) require no interpolation.

Fix now \(K>0\), set \(\mu=\rho K\), and define
\begin{equation}
  A_{\rho,K}(z)=G_K(z)-G_\mu(z).
  \label{eq:sc-shell-action}
\end{equation}
For \(0\le z\le K\), let
\begin{equation}
  s_\mu(z)=
  \begin{cases}
  \min\{H_\mu(z),1/4\},&0\le z<\mu,\\
  1/4,&\mu\le z\le K,
  \end{cases}
  \qquad
  U_{\rho,K}(z)=A_{\rho,K}(z)+s_\mu(z).
  \label{eq:sc-optimized-proxy}
\end{equation}
The second line, rather than \(H_\mu\), is always used at \(z=\mu\).

\begin{proposition}[Strict phase-to-lattice reduction]
\label{prop:sc-phase-lattice}
For \(0<\rho<1\) and \(K>0\),
\begin{equation}
  \gamma_{\rho,K}(z)<U_{\rho,K}(z)
  \le A_{\rho,K}(z)+\frac14,
  \qquad 0\le z\le K.
  \label{eq:sc-pointwise-phase-proxy}
\end{equation}
Because every \(G_a\) is zero-extended, \(A_{\rho,K}\) is the piecewise
action \(G_K-G_{\rho K}\) below the inner interface and \(G_K\) at and
above it.  Consequently,
\begin{align}
  N_D(A_{\rho,1},K^2)
  &\le
  \sum_{\nu_\ell<K}2\nu_\ell
  [U_{\rho,K}(\nu_\ell)]_<
  \label{eq:sc-optimized-spectral-proxy}\\
  &\le
  S_{\rho,K}:=
  \sum_{\nu_\ell<K}2\nu_\ell
  \left\lfloor A_{\rho,K}(\nu_\ell)+\frac14\right\rfloor,
  \qquad \nu_\ell=\ell+\frac12.
  \label{eq:sc-coarse-spectral-proxy}
\end{align}
The cutoff is strict: an order \(\nu_\ell=K\) is inactive.  At a true phase
wall the radial count is one less than the ordinary floor; at a proxy wall
the ordinary floor in \eqref{eq:sc-coarse-spectral-proxy} is retained.
Moreover,
\begin{equation}
  \int_0^\infty2zA_{\rho,K}(z)\,dz
  =\frac{2}{9\pi}(1-\rho^3)K^3.
  \label{eq:sc-exact-action-integral}
\end{equation}
\end{proposition}

\begin{proof}
For \(0\le z<\mu\), the first upper bound in
\eqref{eq:sc-global-difference-enclosure} reads
\[
  \gamma_{\rho,K}(z)
  <A_{\rho,K}(z)+G_\mu(z)-\mathcal F_\mu(z).
\]
From \eqref{eq:sc-auxiliary-F},
\[
  G_\mu(z)-\mathcal F_\mu(z)
  =\min\{H_\mu(z),G_\mu(z)+1/4\}.
\]
The upper half of \eqref{eq:sc-action-difference-enclosure} also gives
\(\gamma_{\rho,K}(z)<A_{\rho,K}(z)+1/4\).  Taking the better of these two
strict upper bounds, and using \(G_\mu(z)+1/4\ge1/4\), gives
\[
  \gamma_{\rho,K}(z)
  <A_{\rho,K}(z)+\min\{H_\mu(z),1/4\}
  =U_{\rho,K}(z).
\]
For \(\mu\le z\le K\), one has \(G_\mu(z)=0\), and
\eqref{eq:sc-global-difference-enclosure} gives
\[
  \gamma_{\rho,K}(z)<G_K(z)+1/4=U_{\rho,K}(z).
\]

Orders \(z\ge K\) do not contribute.  Indeed,
\eqref{eq:sc-individual-phase-enclosure} gives
\[
  \frac{\theta_z(K)}\pi<-\frac14,
  \qquad
  \frac{\theta_z(\mu)}\pi>-\frac12,
\]
because \(G_K(z)=0\) and \(\mathcal F_\mu(z)=-1/4\).  Together with phase
monotonicity this yields
\[
  0<\gamma_{\rho,K}(z)<\frac14<1,
\]
so there is no positive phase level below \(K\).  This includes \(z=K\).

For positive \(x<y\), every positive integer strictly below \(x\) is also
strictly below \(y\), so \([x]_<\le[y]_<\le\lfloor y\rfloor\).  Applying
this observation to the exact phase count
\eqref{eq:sc-exact-phase-count}, and then using
\(s_\mu\le1/4\), proves
\eqref{eq:sc-optimized-spectral-proxy}--
\eqref{eq:sc-coarse-spectral-proxy} without changing either wall
convention.

Finally, scaling \(z=ax\) gives
\[
  \int_0^a2zG_a(z)\,dz
  =\frac{2a^3}{\pi}
  \left(\int_0^1x\sqrt{1-x^2}\,dx
  -\int_0^1x^2\arccos x\,dx\right).
\]
The two elementary integrals are \(1/3\) and \(2/9\), respectively; hence
the last expression is \(2a^3/(9\pi)\).  Subtracting the result with
\(a=\mu=\rho K\) from that with \(a=K\) proves
\eqref{eq:sc-exact-action-integral}.
\end{proof}

The necessity of retaining integerization can be displayed exactly.  Put
\[
  m_K=\max\{0,\lceil K-\tfrac12\rceil\},
  \qquad
  r_\mu=\max\{0,\lceil\mu-\tfrac12\rceil\}.
\]
Thus the active orders are \(\nu_0,\ldots,\nu_{m_K-1}\), and precisely the
first \(r_\mu\) of them lie strictly below \(\mu\).  Define
\[
  E_0(a)=\sum_{\nu_\ell<a}2\nu_\ell G_a(\nu_\ell)
  -\frac{2a^3}{9\pi},
  \qquad
  E_{\mathrm{mp}}(\rho,K)=E_0(K)-E_0(\mu),
\]
\[
  E_{\mathrm{ph}}
  =\sum_{\ell=0}^{r_\mu-1}(2\ell+1)
  \min\{H_\mu(\nu_\ell),1/4\}
  +\frac{m_K^2-r_\mu^2}{4}.
\]
For \(x>0\), define the strict remainder
\[
  \langle x\rangle_<:=x-[x]_<\in(0,1],
\]
so \(\langle m\rangle_<=1\) at positive integer \(m\), and set
\[
  E_{\mathrm{frac}}
  =\sum_{\nu_\ell<K}2\nu_\ell
  \langle U_{\rho,K}(\nu_\ell)\rangle_<.
\]
Expanding \(U=A+s_\mu\) term by term and using
\eqref{eq:sc-exact-action-integral} gives the exact identity
\begin{equation}
  \sum_{\nu_\ell<K}2\nu_\ell[U_{\rho,K}(\nu_\ell)]_<
  -\frac{2}{9\pi}(1-\rho^3)K^3
  =E_{\mathrm{mp}}+E_{\mathrm{ph}}-E_{\mathrm{frac}}.
  \label{eq:sc-fractional-budget}
\end{equation}
In particular, the blanket quarter-slack cost is exactly
\[
  \frac14\sum_{\ell=0}^{m_K-1}(2\ell+1)=\frac{m_K^2}{4};
\]
it cannot be discarded merely because it is of lower polynomial order for
fixed \(\rho\).

% The former small-hole shifted-tail route is retained in source for audit
% history but is not part of the revised proof.
\iffalse
\subsection{Multiplicity reduction to shifted one-dimensional tails}

We record the two external floor-sum inequalities used in the small-hole
argument.  For integers \(a<b\), write
\begin{equation}
  \mathbf T(h;a,b)
  :=\frac12\lfloor h(a)\rfloor
  +\sum_{j=a+1}^{b-1}\lfloor h(j)\rfloor
  +\frac12\lfloor h(b)\rfloor.
  \label{eq:sc-trapezoidal-floor}
\end{equation}

\begin{lemma}[External trapezoidal floor bounds]
\label{lem:sc-external-floor-bounds}
The following statements hold.
\begin{enumerate}
  \item If \(h\) is nonnegative and concave on an integer interval
  \([a,b]\), then
  \begin{equation}
    \mathbf T(h;a,b)\le\int_a^bh(t)\,dt.
    \label{eq:sc-concave-floor-bound}
  \end{equation}
  \item Suppose \(h\) is nonnegative, decreasing, concave, and
  \(c\)-Lipschitz on an
  integer interval \([a,b]\), where \(0<c<1\).  If an integer
  \(p\in[a,b)\) is the last point of the first ordinary-floor shelf, so that
  \[
    \lfloor h(a)\rfloor=\lfloor h(p)\rfloor
    >\lfloor h(p+1)\rfloor,
  \]
  then
  \begin{equation}
    \mathbf T(h;a,b)
    \leq\int_a^b h(t)\,dt-
          \frac{1-c}{2}(b-p).
    \label{eq:sc-improved-concave-floor-bound}
  \end{equation}
  If no such drop occurs, set \(p=b\); the basic bound in item~1 gives
  the same displayed formula with zero correction.  This last no-drop case
  is a consequence of the basic proposition, not an invocation of the
  improved theorem at \(p=b\).
  \item If \(g\) is nonnegative, decreasing, convex, and
  \(1/2\)-Lipschitz on an integer interval \([a,b]\), if \(g(b)=0\), and if
  \(g\) is \(1/3\)-Lipschitz on \([t_*,b]\) for some
  \(t_*\in[a,b]\), then
  \begin{equation}
    \mathbf T(g+1/4;a,b)
    \le\int_a^bg(t)\,dt-\frac14\lfloor g(t_*)\rfloor.
    \label{eq:sc-improved-convex-floor-bound}
  \end{equation}
  \item If \(g\) is nonnegative, decreasing, convex, and
  \(1/2\)-Lipschitz on \([0,b]\), where \(b>0\) is arbitrary and \(g(b)=0\),
  then
  \begin{equation}
    \left\lfloor g(0)+\frac14\right\rfloor
    +2\sum_{j=1}^{\lfloor b\rfloor}
    \left\lfloor g(j)+\frac14\right\rfloor
    \le2\int_0^bg(t)\,dt.
    \label{eq:sc-convex-shifted-tail-bound}
  \end{equation}
\end{enumerate}
\end{lemma}

The first assertion is Proposition~3.1 of \cite{FLPSannuli}; the
first-shelf improvement is its Theorem~1.4; and the third assertion is its
Theorem~3.4.  The fourth is the translated convex shifted-floor theorem of
\cite{FLPSballs}.  These published one-dimensional inequalities are the only
non-special-function external inputs in what follows.

For \(r\in\mathbb N_0\), put \(x_r=r+\tfrac12\) and define the coarse shifted
tail
\begin{equation}
  \mathcal T_r=
  \left\lfloor A_{\rho,K}(x_r)+\frac14\right\rfloor
  +2\sum_{j\ge1}
  \left\lfloor A_{\rho,K}(x_r+j)+\frac14\right\rfloor.
  \label{eq:sc-shifted-tail}
\end{equation}
All but finitely many terms vanish.

\begin{lemma}[Weighted-lattice scaffold]
\label{lem:sc-weighted-scaffold}
One has
\begin{equation}
  S_{\rho,K}=\sum_{r\ge0}\mathcal T_r.
  \label{eq:sc-tail-decomposition}
\end{equation}
If, for every \(r\) with \(x_r<K\),
\begin{equation}
  \mathcal T_r\le2\int_{x_r}^{K}A_{\rho,K}(z)\,dz,
  \label{eq:sc-tail-target}
\end{equation}
then
\begin{equation}
  S_{\rho,K}\le
  \frac{2}{9\pi}(1-\rho^3)K^3.
  \label{eq:sc-scaffold-conclusion}
\end{equation}
More generally, the same conclusion holds under the single aggregate
hypothesis
\begin{equation}
 \sum_{r\geq0}\left(
 \mathcal T_r-2\int_{x_r}^{K}A_{\rho,K}(z)\,dz
 \right)\leq0.
 \label{eq:sc-aggregate-tail-target}
\end{equation}
Thus the per-tail condition \eqref{eq:sc-tail-target} is a sufficient
special case, not an additional requirement on arguments that establish
aggregate cancellation directly.
\end{lemma}

\begin{proof}
For any finitely supported nonnegative integer sequence \(q_\ell\),
\[
  \sum_{\ell\ge0}(2\ell+1)q_\ell
  =\sum_{r\ge0}\left(q_r+2\sum_{\ell>r}q_\ell\right).
\]
Indeed, \(q_\ell\) appears once with coefficient one when \(r=\ell\) and
appears \(\ell\) times with coefficient two when \(0\le r<\ell\).  Taking
\[
  q_\ell=\left\lfloor
  A_{\rho,K}(\ell+\tfrac12)+\frac14\right\rfloor
\]
proves \eqref{eq:sc-tail-decomposition}; samples at or beyond \(K\) are
zero.

Either \eqref{eq:sc-tail-target} or
\eqref{eq:sc-aggregate-tail-target}, together with
\eqref{eq:sc-tail-decomposition}, gives
\[
  S_{\rho,K}
  \le2\int_0^K f_3(z)A_{\rho,K}(z)\,dz,
  \qquad
  f_3(z)=\#\{r\in\mathbb N_0:r+\tfrac12\le z\}.
\]
Let \(F_3(z)=\int_0^zf_3(t)\,dt\).  If
\(z=m+\tfrac12+s\), where \(m\in\mathbb N_0\) and \(0\le s<1\), then direct
summation gives
\begin{equation}
  \frac{z^2}{2}-F_3(z)
  =\frac12\left(s-\frac12\right)^2\ge0;
  \label{eq:sc-step-primitive}
\end{equation}
for \(0\le z<1/2\) the same inequality is immediate because \(F_3(z)=0\).
The action \(A_{\rho,K}\) is absolutely continuous, decreasing, and vanishes
at \(K\).  Indeed, away from the inner interface,
\[
  A_{\rho,K}'(z)=
  \begin{cases}
  \displaystyle-\frac1\pi\left(
  \arccos\frac zK-\arccos\frac z\mu\right)<0,
  &0<z<\mu,\\[2mm]
  \displaystyle-\frac1\pi\arccos\frac zK<0,
  &\mu<z<K,
  \end{cases}
\]
and the two formulas have the same limiting value at \(z=\mu\).
Integration by parts and \eqref{eq:sc-step-primitive} therefore
give
\[
  2\int_0^Kf_3A_{\rho,K}
  =-2\int_0^KF_3A_{\rho,K}'
  \le-\int_0^Kz^2A_{\rho,K}'(z)\,dz
  =\int_0^K2zA_{\rho,K}(z)\,dz.
\]
Now apply \eqref{eq:sc-exact-action-integral}.
\end{proof}

\begin{lemma}[Tails at or above the inner interface]
\label{lem:sc-high-interface-tails}
If \(x_r\ge\mu=\rho K\) and \(x_r<K\), then
\[
  \mathcal T_r\le2\int_{x_r}^{K}A_{\rho,K}(z)\,dz.
\]
Equality \(x_r=\mu\) is included.
\end{lemma}

\begin{proof}
On \(z\ge x_r\ge\mu\), one has \(G_\mu(z)=0\), including at \(z=\mu\),
so \(A_{\rho,K}(z)=G_K(z)\).  Set \(b=K-x_r>0\) and
\(g(t)=G_K(x_r+t)\) on \([0,b]\).  Direct differentiation of
\eqref{eq:sc-model-action} gives
\[
  G_K'(z)=-\frac1\pi\arccos\frac zK\in[-1/2,0],
  \qquad
  G_K''(z)=\frac1{\pi\sqrt{K^2-z^2}}>0.
\]
Thus \(g\) is nonnegative, decreasing, convex, and \(1/2\)-Lipschitz, with
\(g(b)=G_K(K)=0\).  Terms of \eqref{eq:sc-shifted-tail} with \(j>b\) vanish;
if \(b\) is an integer, the terminal term is
\(\lfloor G_K(K)+1/4\rfloor=0\).  Therefore
\eqref{eq:sc-convex-shifted-tail-bound} is exactly the desired estimate.
\end{proof}

\subsection{The low-interface estimate in the small-hole sector}

Define
\begin{equation}
  \omega_0:=G_1(1/2)=\frac{\sqrt3}{2\pi}-\frac16,
  \qquad
  C_*:=\frac12+\frac{8\sqrt2}{15}.
  \label{eq:sc-small-hole-constants}
\end{equation}
Notice that \(0<\omega_0<1/2\).  The lower inequality is equivalent to
\(\pi<3\sqrt3\), which follows from \(\pi<4<3\sqrt3\); the upper inequality
also follows directly from \(0<G_1(1/2)<G_1(0)=1/\pi<1/2\).

\begin{lemma}[Low-interface shifted tails]
\label{lem:sc-low-interface-tails}
Assume
\begin{equation}
  0<\rho<\omega_0,
  \qquad
  K(\omega_0-\rho)\ge C_*.
  \label{eq:sc-low-interface-sector}
\end{equation}
If \(x_r=r+\tfrac12<\mu=\rho K\), then, in fact strictly,
\begin{equation}
  \mathcal T_r
  <2\int_{x_r}^{K}A_{\rho,K}(z)\,dz.
  \label{eq:sc-low-tail-conclusion}
\end{equation}
\end{lemma}

\begin{proof}
Fix a low tail and abbreviate \(x_0=x_r\), \(A=A_{\rho,K}\).  Put
\begin{equation}
  n=\lfloor\mu-x_0\rfloor,
  \qquad q=x_0+n,
  \qquad B=\lceil K-x_0\rceil.
  \label{eq:sc-interface-bookkeeping}
\end{equation}
Then
\[
  q\le\mu<q+1,
  \qquad n<B,
  \qquad G_K(x_0+B)=0.
\]
The point \(q\) is the largest half-integer in the tail grid that does not
exceed \(\mu\); it need not equal \(\lfloor\mu\rfloor\).  A low tail also
forces \(\mu>x_0\ge1/2\), hence \(\mu>1/2\).

Let
\[
  h_A(t)=A(x_0+t)+\frac14,
  \qquad h_K(t)=G_K(x_0+t)+\frac14.
\]
All samples at and beyond \(B\) have floor zero, so
\[
  \frac{\mathcal T_r}{2}=\mathbf T(h_A;0,B).
\]
If \(n\ge1\), additivity of the trapezoidal sum at the shared endpoint and
\(A\le G_K\) give
\begin{equation}
  \frac{\mathcal T_r}{2}
  \le\mathbf T(h_A;0,n)+\mathbf T(h_K;n,B).
  \label{eq:sc-tail-split-positive-n}
\end{equation}
At the split point, the two half-weights add to the original full sample;
after replacing the suffix by \(G_K\), its half-sample can only increase.
If \(n=0\), there is no concave block, and instead
\begin{equation}
  \frac{\mathcal T_r}{2}\le\mathbf T(h_K;0,B).
  \label{eq:sc-tail-split-zero-n}
\end{equation}

For \(0<z<\mu\),
\begin{equation}
  A''(z)=\frac1{\pi\sqrt{K^2-z^2}}
  -\frac1{\pi\sqrt{\mu^2-z^2}}<0.
  \label{eq:sc-action-concavity}
\end{equation}
Thus \(A\) is concave on \([x_0,q]\), also when \(q=\mu\) by continuity.
For \(n\ge1\), the concave bound
\eqref{eq:sc-concave-floor-bound} gives
\begin{equation}
  \mathbf T(h_A;0,n)
  \le\int_0^nA(x_0+t)\,dt+\frac n4.
  \label{eq:sc-concave-head}
\end{equation}

For the suffix set \(g(t)=G_K(x_0+t)\), using the zero extension after
\(K\).  It is nonnegative, decreasing, convex, globally
\(1/2\)-Lipschitz, and \(g(B)=0\).  Because
\(q\le\mu=\rho K<\omega_0K<K/2\),
\[
  t_*:=\frac K2-x_0
\]
belongs to \([n,B]\) (strictly to the right of \(n\)).  On
\([t_*,B]\),
\[
  |g'(t)|\le\frac1\pi\arccos\frac12=\frac13,
\]
and
\[
  g(t_*)=G_K(K/2)=\omega_0K.
\]
Every hypothesis of \eqref{eq:sc-improved-convex-floor-bound} is therefore
satisfied, and
\begin{equation}
  \mathbf T(h_K;n,B)
  \le\int_n^BG_K(x_0+t)\,dt
  -\frac14\lfloor\omega_0K\rfloor.
  \label{eq:sc-convex-suffix}
\end{equation}
The identical statement with \(n=0\) applies to
\eqref{eq:sc-tail-split-zero-n}.

Combining \eqref{eq:sc-concave-head} and
\eqref{eq:sc-convex-suffix}, or using only the latter when \(n=0\), leaves
one and only one integral mismatch:
\begin{equation}
  \delta:=\int_q^\mu G_\mu(z)\,dz.
  \label{eq:sc-interface-loss}
\end{equation}
Indeed,
\[
  \int_{x_0}^{q}A(z)\,dz+\int_q^KG_K(z)\,dz
  =\int_{x_0}^KA(z)\,dz+\delta,
\]
and the same identity holds for \(n=0\), where \(q=x_0\).  Hence
\begin{equation}
  \frac{\mathcal T_r}{2}
  \le\int_{x_0}^KA(z)\,dz
  +\delta-\frac{\lfloor\omega_0K\rfloor-n}{4}.
  \label{eq:sc-tail-compensation}
\end{equation}

Lemma~\ref{lem:sc-turning-bound} gives
\(G_\mu(z)<(\mu-z)^{3/2}/(3\sqrt\mu)\) for
\(0\leq z<\mu\).
If \(q<\mu\), integration gives a strict first bound; if \(q=\mu\), both
sides below vanish.  The endpoint-safe statement is
\begin{equation}
  0\le\delta
  \le\frac{2(\mu-q)^{5/2}}{15\sqrt\mu}
  <\frac{2\sqrt2}{15},
  \label{eq:sc-interface-loss-bound}
\end{equation}
because \(0\le\mu-q<1\) and \(\mu>1/2\).  The last inequality remains
strict when \(q=\mu\), since then \(\delta=0\).

Finally, \(x_0\ge1/2\) implies
\[
  n\le\mu-x_0\le\rho K-\frac12,
\]
while \(\lfloor y\rfloor>y-1\) for every real \(y\).  Therefore
\begin{align*}
  \lfloor\omega_0K\rfloor-n
  &>\omega_0K-1-\left(\rho K-\frac12\right)\\
  &=K(\omega_0-\rho)-\frac12\\
  &\ge\frac{8\sqrt2}{15}
  >4\delta.
\end{align*}
Both strict inequalities remain strict when the threshold in
\eqref{eq:sc-low-interface-sector} is attained.  Substitution in
\eqref{eq:sc-tail-compensation} proves
\eqref{eq:sc-low-tail-conclusion}.
\end{proof}

\subsection{Uniform closure of the small-hole endpoint}

Set
\begin{equation}
  \rho_*:=\frac{\omega_0}{1+2C_*}
  =\frac{\frac{\sqrt3}{2\pi}-\frac16}
  {2+\frac{16\sqrt2}{15}}.
  \label{eq:sc-rho-star}
\end{equation}
Since \(\omega_0>0\) and \(C_*>0\),
\(0<\rho_*<\omega_0\).  Its defining identity is
\begin{equation}
  \omega_0-\rho_*=2C_*\rho_*,
  \qquad
  \frac{\omega_0-\rho_*}{2\rho_*}=C_*.
  \label{eq:sc-rho-star-identity}
\end{equation}

\begin{theorem}[Complete small-hole endpoint]
\label{thm:sc-small-hole-endpoint}
For every \(0<\rho\le\rho_*\) and every \(K\ge0\),
\begin{equation}
  \boxed{
  N_D(A_{\rho,1},K^2)
  \le\frac{2}{9\pi}(1-\rho^3)K^3.}
  \label{eq:sc-small-hole-polya}
\end{equation}
The face \(\rho=\rho_*\), the interface wall \(\rho K=1/2\), and every
strict spectral wall are included.
\end{theorem}

\begin{proof}
If \(K=0\), the strict spectrum is positive, so the left side and the right
side of \eqref{eq:sc-small-hole-polya} are both zero.  Assume \(K>0\) and
write \(\mu=\rho K\).

First suppose \(\mu\le1/2\).  Every shifted-tail start satisfies
\[
  x_r=r+\frac12\ge\frac12\ge\mu.
\]
Lemma~\ref{lem:sc-high-interface-tails} therefore proves
\eqref{eq:sc-tail-target} for every active tail.  At the joining wall
\(\mu=1/2\), the first start is exactly \(x_0=\mu\), which belongs to the
inclusive high-interface lemma.

Now suppose \(\mu>1/2\).  Then
\begin{equation}
  K>\frac1{2\rho}.
  \label{eq:sc-high-K-from-interface}
\end{equation}
For \(0<\rho\le\rho_*\), the defining identity
\eqref{eq:sc-rho-star-identity} gives
\begin{equation}
  \frac{\omega_0-\rho}{2\rho}\ge C_*;
  \label{eq:sc-uniform-small-hole-threshold}
\end{equation}
equivalently, multiply
\((1+2C_*)\rho\le(1+2C_*)\rho_*=\omega_0\)
by positive quantities.  Since
\(\omega_0-\rho>0\), \eqref{eq:sc-high-K-from-interface} and
\eqref{eq:sc-uniform-small-hole-threshold} imply
\begin{equation}
  K(\omega_0-\rho)
  >\frac{\omega_0-\rho}{2\rho}
  \ge C_*.
  \label{eq:sc-strict-small-hole-threshold}
\end{equation}
Every start \(x_r\ge\mu\) is controlled by
Lemma~\ref{lem:sc-high-interface-tails}; every start \(x_r<\mu\) is
controlled by Lemma~\ref{lem:sc-low-interface-tails}, because
\(\rho\le\rho_*<\omega_0\) and
\eqref{eq:sc-strict-small-hole-threshold} supplies its threshold.  Thus all
active tails satisfy \eqref{eq:sc-tail-target} in this branch as well.

Lemma~\ref{lem:sc-weighted-scaffold} now yields
\[
  S_{\rho,K}\le\frac{2}{9\pi}(1-\rho^3)K^3,
\]
and Proposition~\ref{prop:sc-phase-lattice} gives
\(N_D(A_{\rho,1},K^2)\le S_{\rho,K}\), proving
\eqref{eq:sc-small-hole-polya}.

At \(\rho=\rho_*\), inequality
\eqref{eq:sc-uniform-small-hole-threshold} is an equality, but the first
inequality in \eqref{eq:sc-strict-small-hole-threshold} is strict whenever
\(\mu>1/2\).  Hence the endpoint is included.  If instead
\(K(\omega_0-\rho)=C_*\), then
\[
  \mu=\frac{C_*\rho}{\omega_0-\rho}\le\frac12
\]
by \eqref{eq:sc-uniform-small-hole-threshold}, so the point already belongs
to the all-high-interface branch.  No limiting argument at \(\rho=0\), no
continuity of Bessel roots, and no finite numerical verification is used.
\end{proof}

\fi

\section{Product and complementary-action tools}
\label{sec:product-action-tools}

We now prove the P\'olya inequality uniformly as the shell becomes thin.  In
this section the counting function is strict:
\[
  N_D(\Omega,K^2)=\#\{j:\lambda_j(\Omega)<K^2\}.
\]
We retain the global convention
\([u]_<:=\#\{n\in\mathbb N:n<u\}=\max\{0,\lceil u\rceil-1\}\).
Thus a spectral level lying exactly on the wall \(K^2\) is not counted.  Put
\begin{equation}
  \varepsilon=1-\rho,
  \qquad a=\varepsilon K,
  \qquad
  m_\varepsilon=1-\varepsilon+\frac{\varepsilon^2}{3}.
  \label{thin:eq-normalized-variables}
\end{equation}
Since \(1-\rho^3=3\varepsilon m_\varepsilon\), the desired right-hand side
has the two equivalent forms
\begin{equation}
  \frac{2}{9\pi}(1-\rho^3)K^3
  =\frac{1}{\varepsilon^2}\frac{2a^3}{3\pi}m_\varepsilon.
  \label{thin:eq-weyl-normalization}
\end{equation}

The proof uses the optimized phase estimate already established in
Proposition~\ref{prop:sc-phase-lattice}.  We record precisely the consequence
that is used here.  For
\(\nu_\ell=\ell+\tfrac12\), let
\(\gamma_{\rho,K}(\nu_\ell)\) be the normalized radial Bessel phase.  If
\(0\leq\nu_\ell\leq K\), that proposition gives
\begin{equation}
  \gamma_{\rho,K}(\nu_\ell)
  < A_{\rho,K}(\nu_\ell)+\frac14,
  \label{thin:eq-phase-input}
\end{equation}
where, after the scaling in \eqref{thin:eq-normalized-variables},
\(A_{\rho,K}(\nu)=\mathcal A(\varepsilon\nu)\) and
\(\mathcal A\) is the explicit action in
\eqref{thin:eq-action-definition} below.  The exact radial count in channel
\(\ell\) is
\begin{equation}
  \#\{n\geq1:n<\gamma_{\rho,K}(\nu_\ell)\}.
  \label{thin:eq-exact-phase-count}
\end{equation}
In particular, \eqref{thin:eq-phase-input} is an inequality for the strict
count; it is not a replacement of that count by an ordinary floor.

\subsection{The product comparison}

Under the standard unitary radial transformation, the angular channel
\(\ell\in\mathbb N_0\) is the Dirichlet operator
\[
  L_\ell=-\frac{d^2}{dr^2}+\frac{\ell(\ell+1)}{r^2}
  \quad\hbox{in }L^2((\rho,1),dr),
\]
and it occurs with multiplicity \(2\ell+1\).  The following lemma includes
the direction of the min--max comparison and every strict radial and angular
wall.

\begin{lemma}[Strict product majorant]
\label{thin:lem-product-majorant}
Let \(0<\varepsilon<1\), \(a=\varepsilon K\), and define
\begin{equation}
  N(a)=\max\left\{0,\left\lceil\frac a\pi\right\rceil-1\right\}.
  \label{thin:eq-radial-count}
\end{equation}
For \(1\leq n\leq N(a)\), put
\[
  b_n=\sqrt{a^2-n^2\pi^2},
  \qquad
  M_n=
  \left\lceil
  \sqrt{\left(\frac{b_n}{\varepsilon}\right)^2+\frac14}
  -\frac12
  \right\rceil.
\]
Then
\begin{equation}
  N_D(A_{\rho,1},K^2)
  \leq \mathcal P_\varepsilon(a)
  :=\sum_{n=1}^{N(a)}M_n^2.
  \label{thin:eq-product-majorant}
\end{equation}
In particular, the strict shell count is zero for \(0\leq a\leq\pi\).
\end{lemma}

\begin{proof}
The two quadratic forms have the common domain \(H_0^1(\rho,1)\), and
\[
  \int_\rho^1\left(
     |u'|^2+\frac{\ell(\ell+1)}{r^2}|u|^2
  \right)dr
  \geq
  \int_\rho^1\left(|u'|^2+\ell(\ell+1)|u|^2\right)dr,
\]
because \(r\leq1\).  The min--max principle therefore gives
\[
  \lambda_{\ell,n}
  \geq \left(\frac{n\pi}{\varepsilon}\right)^2
       +\ell(\ell+1),
  \qquad n\geq1.
\]
Consequently, a shell eigenvalue in this channel can be counted below
\(K^2\) only if
\begin{equation}
  n^2\pi^2+\varepsilon^2\ell(\ell+1)<a^2.
  \label{thin:eq-product-condition}
\end{equation}
The active radial integers are exactly those satisfying \(n\pi<a\), whose
number is \eqref{thin:eq-radial-count}.  For a fixed active \(n\), condition
\eqref{thin:eq-product-condition} is
\[
  \ell<
  \sqrt{\left(\frac{b_n}{\varepsilon}\right)^2+\frac14}-\frac12.
\]
Hence the allowed angular indices are
\(0,\ldots,M_n-1\), and their full spherical multiplicity is
\[
  \sum_{\ell=0}^{M_n-1}(2\ell+1)=M_n^2.
\]
Summing proves \eqref{thin:eq-product-majorant}.

If \(a=m\pi\), the radial level \(n=m\) is excluded by the strict
inequality \(n\pi<a\); thus \(N(a)=m-1\).  At an angular wall
\(b_n^2/\varepsilon^2=L(L+1)\), the displayed expression inside the ceiling
is exactly \(L\), so \(M_n=L\) and the level \(\ell=L\) is excluded.  The
jump of multiplicity \(2L+1\) occurs only above the wall.  Finally, no
positive integer satisfies \(n\pi<a\) when \(0\leq a\leq\pi\), including
both endpoints, so the product majorant and hence the shell count are zero
there.
\end{proof}


\subsection{The complementary-action comparison}

For \(r\in\{\rho,1\}\) and \(0\leq y\leq ar\), define
\begin{equation}
  J_r(y)=\sqrt{a^2r^2-y^2}
          -y\arccos\frac{y}{ar},
  \label{thin:eq-J-action}
\end{equation}
and set
\begin{equation}
  \mathcal A(y)=
  \begin{cases}
  \dfrac{J_1(y)-J_\rho(y)}{\pi\varepsilon},
    &0\leq y\leq\rho a,\\[6pt]
  \dfrac{J_1(y)}{\pi\varepsilon},
    &\rho a\leq y\leq a.
  \end{cases}
  \label{thin:eq-action-definition}
\end{equation}
We extend \(\mathcal A(y)=0\) for \(y>a\).

\begin{lemma}[Action geometry and curvature switch]
\label{thin:lem-action-geometry}
The function \(\mathcal A\) is continuous, continuously differentiable at
\(y=\rho a\), and strictly decreasing from
\[
  T:=\mathcal A(0)=\frac a\pi
  \quad\hbox{to}\quad \mathcal A(a)=0.
\]
Let \(R:[0,T]\to[0,a]\) be its decreasing inverse, let
\[
  \tau=\mathcal A(\rho a),
  \qquad F(t)=R(t)^2,
  \qquad g(t)=-F'(t).
\]
Then
\begin{equation}
\begin{gathered}
  g\text{ decreases on }(0,\tau),
  \qquad g\text{ increases on }(\tau,T),\\
  F(0)=a^2,
  \qquad F(\tau)=\rho^2a^2,
  \qquad F(T)=0,
  \qquad g(T)=2\pi\rho a,
\end{gathered}
  \label{thin:eq-action-geometry}
\end{equation}
where \(g(T)\) denotes the one-sided limit.  Moreover \(g(t)=O(t^{-1/3})\)
as \(t\downarrow0\).
\end{lemma}

\begin{proof}
Direct differentiation on the appropriate open intervals gives
\begin{equation}
  J_r'(y)=-\arccos\frac{y}{ar},
  \qquad
  J_r''(y)=\frac1{\sqrt{a^2r^2-y^2}}.
  \label{thin:eq-J-derivatives}
\end{equation}
At \(y=\rho a\), \(J_\rho(\rho a)=0\), and the two one-sided first
derivatives of \eqref{thin:eq-action-definition} both equal
\(-\arccos\rho/(\pi\varepsilon)\).  Thus the two pieces and their first
derivatives match.  Both derivatives are negative away from \(y=a\), so the
inverse exists on exactly the stated interval.

On the outer branch \(\rho a<y<a\), put
\(\alpha=\arccos(y/a)\).  If \(t=\mathcal A(y)\), then
\begin{equation}
  g(t)=-\frac{2y}{\mathcal A'(y)}
      =\frac{2\pi\varepsilon y}{\alpha}.
  \label{thin:eq-g-outer}
\end{equation}
The function \(y/\alpha\) increases with \(y\), whereas \(y=R(t)\)
decreases with \(t\); hence \(g\) decreases on \(0<t<\tau\).

On the inner branch \(0<y<\rho a\), put
\[
  \delta(y)=\arccos\frac ya-\arccos\frac{y}{\rho a}>0.
\]
Differentiating \(\arccos(y/(ar))\) with respect to \(r\) gives the exact
identity
\begin{equation}
  \frac{\delta(y)}{y}
  =\int_\rho^1\frac{dr}{r\sqrt{a^2r^2-y^2}}.
  \label{thin:eq-delta-integral}
\end{equation}
The integrand increases strictly with \(y\), so \(\delta(y)/y\) increases,
and \(y/\delta(y)\) decreases, with \(y\).  Since
\[
  g(t)=\frac{2\pi\varepsilon y}{\delta(y)}
\]
and \(y=R(t)\) decreases, \(g\) increases on \(\tau<t<T\).  Formula
\eqref{thin:eq-g-outer} and its inner counterpart agree at the interface.

The three values of \(F\) follow from
\(R(0)=a\), \(R(\tau)=\rho a\), and \(R(T)=0\).  As \(y\downarrow0\),
\eqref{thin:eq-delta-integral} gives
\[
  \delta(y)=\frac{\varepsilon}{\rho a}y+O(y^3),
\]
and therefore \(g(T)=2\pi\rho a\).

Finally, on the outer branch write \(y=a\cos\theta\).  Then
\begin{equation}
  t=\frac{a}{\pi\varepsilon}
       (\sin\theta-\theta\cos\theta)
    \sim\frac{a\theta^3}{3\pi\varepsilon},
  \qquad
  g(t)=\frac{2\pi\varepsilon a\cos\theta}{\theta}
       =O(t^{-1/3}).
  \label{thin:eq-improper-asymptotic}
\end{equation}
This proves the asserted improper endpoint behavior.  All square roots and
arccosines in the proof are taken on their stated closed domains; every
division occurs in an open interval where its denominator is positive, and
the endpoint zeros are handled by the limits just computed.
\end{proof}

\begin{lemma}[Strict phase-to-action layer cake]
\label{thin:lem-layer-cake}
For \(\ell\in\mathbb N_0\), set
\begin{equation}
  y_\ell=\varepsilon\left(\ell+\frac12\right),
  \qquad
  q_\ell=\left[\mathcal A(y_\ell)+\frac14\right]_<,
  \qquad
  P_{\mathcal A}=\varepsilon\sum_{\ell\geq0}y_\ell q_\ell.
  \label{thin:eq-phase-proxy}
\end{equation}
Then
\begin{equation}
  N_D(A_{\rho,1},K^2)
  \leq\frac{2P_{\mathcal A}}{\varepsilon^2}.
  \label{thin:eq-phase-transfer}
\end{equation}
For \(n\geq1\), put \(x_n=n-\tfrac14\), and define
\begin{equation}
  M_\varepsilon(x)
  =\#\left\{\ell\in\mathbb N_0:
       \varepsilon\left(\ell+\frac12\right)<x\right\}
  =\max\left\{0,
       \left\lceil\frac{x}{\varepsilon}-\frac12\right\rceil
       \right\}.
  \label{thin:eq-half-integer-count}
\end{equation}
The proxy has the exact finite representation
\begin{equation}
  P_{\mathcal A}
  =\frac{\varepsilon^2}{2}
    \sum_{\substack{n\geq1\\x_n<T}}
       M_\varepsilon(R(x_n))^2.
  \label{thin:eq-discrete-layer-cake}
\end{equation}
Moreover,
\begin{equation}
  I:=\frac12\int_0^T F(t)\,dt
  =\frac{a^3}{3\pi}
     \left(1-\varepsilon+\frac{\varepsilon^2}{3}\right),
  \qquad
  \frac{2I}{\varepsilon^2}
  =\frac{2}{9\pi}(1-\rho^3)K^3.
  \label{thin:eq-continuous-normalization}
\end{equation}
\end{lemma}

\begin{proof}
For \(y_\ell<a\), equations \eqref{thin:eq-phase-input} and
\eqref{thin:eq-exact-phase-count} show that the number of radial levels in
channel \(\ell\) is at most \(q_\ell\).  If \(y_\ell\geq a\), equivalently
\(\nu_\ell\geq K\), the product comparison gives
\[
  \lambda_{\ell,1}
  \geq\frac{\pi^2}{\varepsilon^2}+\ell(\ell+1)
  =\frac{\pi^2}{\varepsilon^2}+\nu_\ell^2-\frac14
  >K^2,
\]
so there is no radial level; at \(y_\ell=a\), the definition also gives
\(q_\ell=[1/4]_< =0\).  Multiplicity \(2\ell+1=2y_\ell/\varepsilon\)
now yields
\[
  N_D(A_{\rho,1},K^2)
  \leq\sum_{\ell\geq0}\frac{2y_\ell}{\varepsilon}q_\ell
  =\frac{2P_{\mathcal A}}{\varepsilon^2},
\]
which proves \eqref{thin:eq-phase-transfer}.  If
\(\mathcal A(y_\ell)+1/4=m\in\mathbb Z\), then \(q_\ell=m-1\), so the
phase-proxy wall is excluded exactly as required.

For every \(y\geq0\), the strict bracket is
\[
  \left[\mathcal A(y)+\frac14\right]_<
  =\#\{n\geq1:x_n<\mathcal A(y)\}.
\]
The sums are finite.  Since \(\mathcal A\) and \(R\) are decreasing,
\(x_n<\mathcal A(y_\ell)\) is equivalent to \(y_\ell<R(x_n)\).  Therefore
\[
\begin{aligned}
  P_{\mathcal A}
  &=\varepsilon
    \sum_{\substack{n\geq1\\x_n<T}}
      \sum_{y_\ell<R(x_n)}y_\ell \\
  &=\varepsilon
    \sum_{\substack{n\geq1\\x_n<T}}
      \varepsilon\sum_{\ell=0}^{M_\varepsilon(R(x_n))-1}
        \left(\ell+\frac12\right) \\
  &=\frac{\varepsilon^2}{2}
    \sum_{\substack{n\geq1\\x_n<T}}
      M_\varepsilon(R(x_n))^2,
\end{aligned}
\]
because \(\sum_{\ell=0}^{M-1}(\ell+1/2)=M^2/2\).  This proves
\eqref{thin:eq-discrete-layer-cake}.  Formula
\eqref{thin:eq-half-integer-count} follows from the same strict counting;
at \(x/\varepsilon=m+1/2\), it equals \(m\), excluding \(\ell=m\).

Differentiation of \eqref{thin:eq-J-action} with respect to \(r\) gives
\[
  \frac{\partial J_r(y)}{\partial r}
  =\sqrt{a^2-\frac{y^2}{r^2}}
  \quad(0\leq y\leq ar).
\]
Consequently the two cases of \eqref{thin:eq-action-definition} combine as
\begin{equation}
  \mathcal A(y)
  =\frac1{\pi\varepsilon}
    \int_\rho^1
      \left(a^2-\frac{y^2}{r^2}\right)_+^{1/2}\,dr,
  \label{thin:eq-action-integral}
\end{equation}
where \(s_+=\max\{s,0\}\).
The area identity for a decreasing inverse and Tonelli's theorem give
\[
  \frac12\int_0^T R(t)^2\,dt
  =\int_0^a y\mathcal A(y)\,dy.
\]
Using \eqref{thin:eq-action-integral} and then \(y=rz\),
\[
\begin{aligned}
  \int_0^a y\mathcal A(y)\,dy
  &=\frac1{\pi\varepsilon}
    \int_\rho^1r^2\,dr
    \int_0^a z\sqrt{a^2-z^2}\,dz \\
  &=\frac{a^3}{3\pi\varepsilon}\int_\rho^1r^2\,dr
   =\frac{a^3}{3\pi}
      \left(1-\varepsilon+\frac{\varepsilon^2}{3}\right).
\end{aligned}
\]
This is \eqref{thin:eq-continuous-normalization}; its final equality follows
from \(1-\rho^3=3\varepsilon m_\varepsilon\) and \(K=a/\varepsilon\).
\end{proof}

\begin{lemma}[Shifted sawtooth and radial quadrature deficit]
\label{thin:lem-radial-deficit}
Let
\[
  C(t)=\#\{n\geq1:n-\tfrac14<t\},
  \qquad w(t)=t-C(t),
  \qquad \mathfrak W(t)=\int_0^t w(s)\,ds,
\]
and put \(\Psi(t)=\mathfrak W(t)-t/4\).  Then, for all \(t\geq0\),
\begin{equation}
  \mathfrak W(t)\geq0,
  \qquad
  -\frac1{32}\leq\Psi(t)\leq\frac3{32}.
  \label{thin:eq-sawtooth-bounds}
\end{equation}
Furthermore, with
\begin{equation}
  D_{\mathrm{rad}}
  =\int_0^T F(t)\,dt
   -\sum_{\substack{n\geq1\\x_n<T}}F(x_n),
  \label{thin:eq-radial-deficit-definition}
\end{equation}
one has
\begin{equation}
  D_{\mathrm{rad}}
  \geq\frac{\rho^2a^2}{4}-\frac{\pi\rho a}{4}.
  \label{thin:eq-radial-deficit-bound}
\end{equation}
\end{lemma}

\begin{proof}
On \(0\leq t\leq3/4\), \(C(t)=0\), and hence
\begin{equation}
  \Psi(t)=\frac12\left(t-\frac14\right)^2-\frac1{32}.
  \label{thin:eq-first-sawtooth-cell}
\end{equation}
For \(x_n=n-1/4\), \(n\geq1\), and \(t=x_n+u\) with
\(0\leq u\leq1\), direct integration gives
\begin{equation}
  \Psi(t)=\frac12\left(u-\frac12\right)^2-\frac1{32}.
  \label{thin:eq-periodic-sawtooth-cell}
\end{equation}
These formulas agree at cell endpoints and prove the two-sided bound for
\(\Psi\).  On the first cell \(\mathfrak W=t^2/2\geq0\).  Thereafter
\(\mathfrak W=t/4+\Psi\geq t/4-1/32>0\), proving
\eqref{thin:eq-sawtooth-bounds}.  At a jump \(x_n\), strict counting gives
the one-sided values \(w(x_n^-)=3/4\) and \(w(x_n^+)=-1/4\), while
\(\mathfrak W\)
and \(\Psi\) are continuous.

Distributionally, \(dw=dt-dC\).  Since \(F(T)=0\), an atom at \(x_n=T\)
has zero weight and is consistent with its exclusion from
\eqref{thin:eq-radial-deficit-definition}.  Stieltjes integration by parts,
first on \([\eta,T]\) and then with \(\eta\downarrow0\), gives
\begin{equation}
  D_{\mathrm{rad}}=\int_0^T w(t)g(t)\,dt.
  \label{thin:eq-stieltjes-start}
\end{equation}
Indeed, \eqref{thin:eq-improper-asymptotic}, together with
\(w(t)=t\) and \(\mathfrak W(t)=t^2/2\) near zero, gives
\(w(t)g(t)=O(t^{2/3})\) and
\(\mathfrak W(t)g(t)=O(t^{5/3})\to0\), so the improper
boundary term vanishes.

Split \eqref{thin:eq-stieltjes-start} at the actual value \(\tau\), with
no assumption that \(\tau\) belongs to any particular sawtooth cell.  On
\((0,\tau)\), \(g\) decreases and \(\mathfrak W\geq0\); hence
\begin{equation}
  \int_0^\tau wg
  =\mathfrak W(\tau)g(\tau)-\int_0^\tau \mathfrak W\,dg
  \geq \mathfrak W(\tau)g(\tau).
  \label{thin:eq-decreasing-stieltjes}
\end{equation}
On \((\tau,T)\), \(g\) increases and \(w=1/4+\Psi'\) almost everywhere.
Adding its integration-by-parts identity to
\eqref{thin:eq-decreasing-stieltjes}, and using
\(\mathfrak W(\tau)=\tau/4+\Psi(\tau)\), cancels the interface value and
yields
\[
  D_{\mathrm{rad}}
  \geq \frac{F(\tau)}4+\frac{\tau g(\tau)}4
     +\Psi(T)g(T)-\int_\tau^T\Psi\,dg.
\]
Because \(dg\) is a positive measure on \((\tau,T)\), the bounds
\eqref{thin:eq-sawtooth-bounds} imply
\[
\begin{aligned}
  D_{\mathrm{rad}}
  &\geq \frac{F(\tau)}4-\frac{g(T)}8
    +g(\tau)\left(\frac\tau4+\frac3{32}\right) \\
  &\geq \frac{F(\tau)}4-\frac{g(T)}8.
\end{aligned}
\]
The two discarded quantities,
\(-\int_0^\tau \mathfrak W\,dg\) and
\(g(\tau)(\tau/4+3/32)\), are explicitly nonnegative.  Substitution of
\(F(\tau)=\rho^2a^2\) and \(g(T)=2\pi\rho a\) from
Lemma~\ref{thin:lem-action-geometry} proves
\eqref{thin:eq-radial-deficit-bound}.  Continuity of
\(\mathfrak W,\Psi,g\) at the
split shows that the same proof covers \(\tau\) below, on, or above every
sawtooth grid point.
\end{proof}


% The former finite compact-middle owner graph is retained in source for
% provenance but is excluded from the live theorem.
\iffalse
\section{The compact-middle argument in full detail}
\label{cm:sec-full}

We retain the strict counting convention
\([x]_<:=\#\{n\in\mathbb N:n<x\}\), put
\[
 W(\rho,K):=\frac{2}{9\pi}(1-\rho^3)K^3,
 \qquad z_\rho:=\frac{\pi}{1-\rho},
 \qquad k_m(\rho):=\sqrt{z_\rho^2+m(m+1)},
\]
and denote, consistently with the separated spectrum above, by
\(k_{\ell,n}(\rho)\) the \(n\)-th positive frequency in
angular channel \(\ell\).  Its multiplicity is \(2\ell+1\).  The phase
reduction proved earlier gives
\begin{equation}
 N_D(A_{\rho,1},K^2)
 \le P_{\rm coarse}(\rho,K)
 :=\sum_{\ell+1/2<K}(2\ell+1)
 \left[A_{\rho,K}(\ell+1/2)+\frac14\right]_<,
 \label{cm:eq-coarse}
\end{equation}
where \(A_{\rho,K}=G_K-G_{\rho K}\) and
\[
 G_a(\nu)=
 \begin{cases}
 \displaystyle\frac1\pi\left(\sqrt{a^2-\nu^2}
 -\nu\arccos\frac\nu a\right),&0\leq\nu<a,\\[4pt]
 0,&\nu\geq a.
 \end{cases}
\]
The value at both interfaces is therefore fixed by zero extension.  This
 will matter at \(\nu=\rho K\), at \(\nu=K\), and at every integer proxy
 wall.
The relation to the ordinary-floor scaffold
\(S_{\rho,K}\) in \eqref{eq:sc-coarse-spectral-proxy} is exact:
\begin{equation}
 P_{\rm coarse}(\rho,K)\leq S_{\rho,K},\qquad
 S_{\rho,K}-P_{\rm coarse}(\rho,K)
 =\sum_{\substack{\ell+1/2<K\\
 A_{\rho,K}(\ell+1/2)+1/4\in\mathbb N}}
 (2\ell+1).
 \label{cm:eq-strict-ordinary-proxy-difference}
\end{equation}
Thus the ordinary floor differs only on positive-integer proxy walls, where
it deliberately keeps the entering layer while the strict bracket does not.

\subsection{Constants and the primitive analytic envelope}

For \(0\leq s\leq1\), set
\[
 G_1(s)=\frac1\pi\left(\sqrt{1-s^2}-s\arccos s\right),
 \quad
 \omega_0=G_1(1/2)=\frac{\sqrt3}{2\pi}-\frac16,
\]
\[
 C_*=\frac12+\frac{8\sqrt2}{15},\qquad
 C_0=1+\frac{8\sqrt2}{15},\qquad
 \rho_*:=\frac{\omega_0}{1+2C_*}=\frac{\omega_0}{2C_0}.
\]
For \(0<\rho<1\), define
\[
 a_\rho=\frac{2\pi\rho}{1-\rho},\qquad
 \eta_\rho=G_1\!\left(\max\{\rho,1/2\}\right),
\]
\begin{equation}
 K_0(\rho)=
 \left(
 \frac{\sqrt{a_\rho}+\sqrt{a_\rho+4\eta_\rho C_0}}
 {2\eta_\rho}
 \right)^2,
 \qquad
 H_0(\rho)=\frac{C_*}{\omega_0-\rho}\quad(\rho<\omega_0).
 \label{cm:eq-upper-functions}
\end{equation}
We first re-establish every primitive owner that is needed to form the
compact residual.  This prevents the residual description from hiding an
unproved coverage assumption.

\begin{proposition}[Primitive analytic owners]
\label{cm:prop-primitive-owners}
The P\'olya bound holds on each of the following inclusive domains:
\begin{enumerate}
\item \(2\rho K\leq1\);
\item \((1-\rho)K\leq\pi\);
\item \(0<\rho<\omega_0\) and
      \(K(\omega_0-\rho)\geq C_*\), equivalently \(K\geq H_0(\rho)\);
\item \(0<\rho\leq\rho_*\), for every \(K\geq0\);
\item \(K\geq K_0(\rho)\);
\end{enumerate}
\end{proposition}

\begin{proof}
Item 1 is exactly Lemma~\ref{lem:sc-high-interface-tails} followed by the
weighted scaffold, Lemma~\ref{lem:sc-weighted-scaffold}.  Item 2 is the
zero-mode conclusion of the strict product majorant,
Lemma~\ref{thin:lem-product-majorant}.  For item 3, the high starts are owned
by Lemma~\ref{lem:sc-high-interface-tails}, the low starts by
Lemma~\ref{lem:sc-low-interface-tails}, and the same scaffold completes the
sum; all three statements include their equality faces.  Item 4 is
Theorem~\ref{thm:sc-small-hole-endpoint}.  These references invoke the full
proofs above and introduce no new convention.

It remains to prove item 5.  Use the shifted starts and tails from
\eqref{eq:sc-shifted-tail}, the trapezoidal sum from
\eqref{eq:sc-trapezoidal-floor}, and the exact weighted decomposition
\eqref{eq:sc-tail-decomposition}.  Fix a low-interface tail
\(x_0=r+1/2<\mu\), and write \(A=A_{\rho,K}\).
Put \(n=\lfloor\mu-x_0\rfloor\), \(q=x_0+n\), and
\(B=\lceil K-x_0\rceil\).  On \([0,n]\),
\(h(t)=A(x_0+t)+1/4\) is decreasing, concave, and Lipschitz with
constant \(c_\rho=\arccos\rho/\pi\).  If \(p\) is the last index of its
first constant ordinary-floor shelf (put \(p=n\) when the shelf never
drops), the concave trapezoidal-floor inequality gives
\begin{equation}
 \mathbf T(h,0,n)
 \leq\int_0^nA(x_0+t)\,dt+\frac n4
 -\frac{1-c_\rho}{2}(n-p).
 \label{cm:eq-fixed-head}
\end{equation}
For \(n=0\), the head is absent and \(p=0\).

Let \(\vartheta=-A'\).  Direct differentiation gives
\[
 \vartheta'(z)=\frac1{\pi\sqrt{\mu^2-z^2}}
 -\frac1{\pi\sqrt{K^2-z^2}}
 \geq\frac{1-\rho}{\pi\rho K}=:\kappa_\rho.
\]
Equal floors on the first shelf imply
\(0\leq A(x_0)-A(x_0+p)<1\), including at an integer endpoint.  Hence
\[
 1>\int_{x_0}^{x_0+p}\vartheta(z)\,dz
 \geq\frac{\kappa_\rho p^2}{2},
 \qquad p<\sqrt{a_\rho K}.
\]

For the convex tail set \(s=\max\{q,K/2\}\).  The function
\(G_K(x_0+t)\) is decreasing and convex, is Lipschitz \(1/2\) on the
whole tail and \(1/3\) after \(s\), and
\(G_K(s)\geq K\eta_\rho\).  The convex trapezoidal-floor inequality
therefore gives
\begin{equation}
 \mathbf T(G_K(x_0+\cdot)+1/4,n,B)
 \leq\int_n^B G_K(x_0+t)\,dt
 -\frac14\lfloor K\eta_\rho\rfloor.
 \label{cm:eq-fixed-tail}
\end{equation}
At \(j=n\), the two half endpoint weights in
\eqref{cm:eq-fixed-head}--\eqref{cm:eq-fixed-tail} majorize the original
full \(A\)-sample.  Their only integral mismatch is
\[
 \delta=\int_q^\mu G_\mu(z)\,dz,
 \qquad 0\leq\delta\leq
 \frac{2(\mu-q)^{5/2}}{15\sqrt\mu}<\frac{2\sqrt2}{15}.
\]
Writing \(d_\rho=1-2c_\rho=2\arcsin\rho/\pi\), combination of the
head and tail yields
\[
 \frac{\mathcal T_r}{2}
 \leq\int_{x_0}^KA(z)\,dz+\delta-\frac M4,
 \quad
 M=\lfloor K\eta_\rho\rfloor+d_\rho(n-p)-p.
\]
Since \(d_\rho>0\) and \(p<\sqrt{a_\rho K}\), the strict floor
inequality, valid even when \(K\eta_\rho\) is an integer, gives
\[
 M>K\eta_\rho-1-\sqrt{a_\rho K}.
\]
Thus the exact sufficient condition is
\begin{equation}
 K\eta_\rho-\sqrt{a_\rho K}\geq C_0.
 \label{cm:eq-fixed-ratio-condition}
\end{equation}
Indeed it implies \(M>8\sqrt2/15>4\delta\), so every low-interface tail
is strictly below twice its action integral.  Tails beginning at or above
\(\mu\) satisfy the non-strict convex high-interface bound; the weighted
identity completes the proof.  The cases \(n=0\), \(p=0\), \(p=n\),
\(q=\mu\), \(q\) on either side of \(K/2\), and all strict-floor walls
are included in the displayed inequalities.

Writing \(Y=\sqrt K\), the left side minus \(C_0\) is
\(\eta_\rho Y^2-\sqrt{a_\rho}Y-C_0\).  It has one positive root, and
the square of that root is precisely \(K_0(\rho)\) in
\eqref{cm:eq-upper-functions}.  Hence \eqref{cm:eq-fixed-ratio-condition}
holds, including equality, exactly for \(K\geq K_0(\rho)\).

\end{proof}

Proposition~\ref{cm:prop-primitive-owners} supplies the inclusive
frequency owners needed in the open ratio strip
\(\rho_*<\rho<39/50\):
\[
 2\rho K\leq1,\qquad (1-\rho)K\leq\pi,\qquad
 \rho<\omega_0,\ K\geq H_0(\rho),\qquad K\geq K_0(\rho).
\]
The endpoint \(\rho=\rho_*\) is already owned by
Theorem~\ref{thm:sc-small-hole-endpoint}; the opposite endpoint
\(\rho=39/50\) will be owned by the all-frequency optical theorem.

Put
\[
 L(\rho)=\max\left\{\frac1{2\rho},\frac\pi{1-\rho}\right\},
 \qquad
 U(\rho)=
 \min\bigl(\{K_0(\rho)\}\cup
 \{H_0(\rho):\rho<\omega_0\}\bigr).
\]
Since every primitive frequency wall is included, exact complementation
within this ratio strip gives
\begin{equation}
 \mathcal D_{16}=
 \left\{(\rho,K):\rho_*<\rho<\frac{39}{50},\quad
 L(\rho)<K<U(\rho)\right\}.
 \label{cm:eq-D16}
\end{equation}
Both frequency inequalities are strict.  No numerical locator for the
switch between \(H_0\) and \(K_0\) is needed; the minimum definition retains
the common face automatically.

The two lower walls meet at
\[
 \rho_c:=\frac1{1+2\pi},\qquad
 L(\rho)=
 \begin{cases}
  (2\rho)^{-1},&\rho\leq\rho_c,\\
  z_\rho,&\rho\geq\rho_c.
 \end{cases}
\]

\subsection{Variational comparisons and the complete zero registry}

The transformed channel form is
\[
 \mathfrak a_\ell[u]=\int_\rho^1
 \left(|u'|^2+\frac{\ell(\ell+1)}{r^2}|u|^2\right)dr,
 \qquad D(\mathfrak a_\ell)=H_0^1(\rho,1).
\]
Since \(r^{-2}\geq1\), the min--max principle gives, at every radial
index,
\begin{equation}
 k_{\ell,n}^2\geq n^2z_\rho^2+\ell(\ell+1),
 \qquad k_{0,n}=nz_\rho.
 \label{cm:eq-product-bound}
\end{equation}
Extension by zero from \((\rho,1)\) to \((0,1)\) embeds the shell
fixed-channel form domain isometrically into the ball fixed-channel form
domain.  The zero trace at \(r=\rho\) preserves membership in
\(H_0^1(0,1)\), and Hardy's inequality controls the centrifugal term.
Thus, preserving both \(\ell\) and \(n\),
\begin{equation}
 k_{\ell,n}^2\geq j_{\ell+1/2,n}^2.
 \label{cm:eq-shell-ball}
\end{equation}
The ball forms have the same domain for every angular order; hence, for
\(p\geq\ell\),
\begin{equation}
 j_{p+1/2,n}^2\geq j_{\ell+1/2,n}^2
 +p(p+1)-\ell(\ell+1).
 \label{cm:eq-angular-propagation}
\end{equation}
These are fixed-channel comparisons, not an invocation of unlabelled
global domain monotonicity.

We use the exact elementary enclosure
\begin{equation}
 \frac{333}{106}<\pi<\frac{22}{7}.
 \label{cm:eq-pi}
\end{equation}
For the lower bound, Machin's identity
\(\pi=16\arctan(1/5)-4\arctan(1/239)\) and alternating sums give a
strictly stronger rational bound; for the upper bound,
\[
 0<\int_0^1\frac{x^4(1-x)^4}{1+x^2}\,dx=\frac{22}{7}-\pi.
\]

For clarity we now derive every specialized Bessel-zero inequality used
below, apart from the one published first-zero estimate that is stated
explicitly before specialization.  The spherical Bessel identity and
recurrences are standard \cite{DLMF}.  Set
\[
 \mathsf j_\ell(x)=\sqrt{\frac{\pi}{2x}}J_{\ell+1/2}(x),
 \qquad
 P_\ell(x)=x^{\ell+1}\mathsf j_\ell(x).
\]
Then
\begin{equation}
 P_{\ell+1}=(2\ell+1)P_\ell-x^2P_{\ell-1},
 \qquad P_\ell'=xP_{\ell-1},
 \label{cm:eq-P-recurrence}
\end{equation}
with \(P_0=\sin x\) and \(P_1=\sin x-x\cos x\).  In particular,
\[
\begin{aligned}
P_2={}&(3-x^2)\sin x-3x\cos x,\\
P_3={}&(15-6x^2)\sin x+(x^3-15x)\cos x,\\
P_5={}&(945-420x^2+15x^4)\sin x
       +(-945x+105x^3-x^5)\cos x,\\
P_6={}&(10395-4725x^2+210x^4-x^6)\sin x\\
 &+(-10395x+1260x^3-21x^5)\cos x,\\
P_7={}&(135135-62370x^2+3150x^4-28x^6)\sin x\\
 &+(-135135x+17325x^3-378x^5+x^7)\cos x.
\end{aligned}
\]
Positive zeros are simple by ODE uniqueness.  Their indices are fixed by
the following standard interlacing argument; the variational angular
propagation \eqref{cm:eq-angular-propagation} is not used for this purpose.

\begin{lemma}[Interlacing determines the zero index from the sign]
\label{cm:lem-zero-index}
For \(\ell\in\mathbb N_0\) and \(n\in\mathbb N\),
\begin{equation}
 j_{\ell+1/2,n}<j_{\ell+3/2,n}<j_{\ell+1/2,n+1}.
 \label{cm:eq-adjacent-interlacing}
\end{equation}
For \(x>0\), set
\[
 Z_\ell(x):=\#\{n\in\mathbb N:j_{\ell+1/2,n}<x\}.
\]
If \(x\) is not a zero of
\(\mathsf j_0,\ldots,\mathsf j_L\), then
\begin{equation}
 Z_{\ell+1}(x)\in
 \bigl\{\max\{Z_\ell(x)-1,0\},Z_\ell(x)\bigr\},
 \qquad 0\leq\ell<L,
 \label{cm:eq-count-two-choices}
\end{equation}
and
\begin{equation}
 \operatorname{sgn}\mathsf j_{\ell+1}(x)
 =(-1)^{Z_{\ell+1}(x)}.
 \label{cm:eq-count-parity}
\end{equation}
Consequently \(Z_0(x)\), together with the signs of the successive
spherical Bessel functions, determines every \(Z_\ell(x)\) uniquely.
\end{lemma}

\begin{proof}
Equation~\eqref{cm:eq-adjacent-interlacing} is exactly the contiguous-order
interlacing formula in DLMF~\cite[Eq.~(10.21.2)]{DLMF}, applied at
\(\nu=\ell+1/2\).  Put \(q=Z_\ell(x)\).  Interlacing shows that the number
of zeros of \(\mathsf j_{\ell+1}\) below \(x\) is at most \(q\) and, when
\(q\geq1\), at least \(q-1\); for \(q=0\) it is zero.  This proves
\eqref{cm:eq-count-two-choices}.  Each \(\mathsf j_\ell\) is positive near
zero and changes sign at every simple positive zero, which proves
\eqref{cm:eq-count-parity}.  The two candidates in
\eqref{cm:eq-count-two-choices} have opposite parity whenever both occur.
\end{proof}

The half-period locations determine \(Z_0(w)\), because
\(j_{1/2,n}=n\pi\).  The recurrence and tangent calculations below give
the displayed nonzero rational-wall sign vectors.  Applying
Lemma~\ref{cm:lem-zero-index} successively from \(\ell=0\) gives the final
column:
\begin{center}
\scriptsize
\renewcommand{\arraystretch}{1.16}
\begin{tabular}{c|c|c|c|c|c}
rational wall \(w\)&target \((\ell,n)\)&location of \(w\)&\(Z_0(w)\)&
signs \(\mathsf j_0,\ldots,\mathsf j_\ell\)&
counts \(Z_0,\ldots,Z_\ell\)\\ \hline
\(77/10\)&\((1,2)\)&\((2\pi,5\pi/2)\)&2&\((+,-)\)&\((2,1)\)\\
\(9\)&\((2,2)\)&\((5\pi/2,3\pi)\)&2&\((+,+,-)\)&\((2,2,1)\)\\
\(10\)&\((6,1)\)&\((3\pi,7\pi/2)\)&3&
\((-,+,+,-,-,-,+)\)&\((3,2,2,1,1,1,0)\)\\
\(103/10\)&\((3,2)\)&\((3\pi,7\pi/2)\)&3&
\((-,+,+,-)\)&\((3,2,2,1)\)\\
\(21/2\)&\((1,3)\)&\((3\pi,7\pi/2)\)&3&\((-,+)\)&\((3,2)\)\\
\(23/2\)&\((7,1)\)&\((7\pi/2,4\pi)\)&3&
\((-,-,+,+,-,-,-,+)\)&\((3,3,2,2,1,1,1,0)\)\\
\(61/5\)&\((2,3)\)&\((7\pi/2,4\pi)\)&3&
\((-,-,+)\)&\((3,3,2)\)\\
\(129/10\)&\((5,2)\)&\((4\pi,9\pi/2)\)&4&
\((+,-,-,+,+,-)\)&\((4,3,3,2,2,1)\)
\end{tabular}
\end{center}
For example, the rational-wall signs follow from \eqref{cm:eq-pi}, the displayed
polynomials, and
\(\tan s>s\), \(\tan s>s+s^3/3\),
\(\tan s<s/(1-s^2/2)\) for \(0<s<1\).  The nonautomatic strict
comparisons can be taken as
\[
\begin{array}{c|c}
77/10&5\pi/2-x>3/20>10/77\\
21/2&7\pi/2-x>49/100>2/21\\
9&3\pi-x>21/50>9/26\\
61/5&4\pi-x>9/25>915/3646\\
103/10&7\pi/2-x>69/100>621540/938227\\
129/10&0<x-4\pi<17/50,\quad\tan(x-4\pi)<1700/4711\\
10&4/7<x-3\pi<29/50,\quad\tan(x-3\pi)<2900/4159\\
23/2&53/50<4\pi-x<\pi/2,\quad
\tan(4\pi-x)>546377/375000.
\end{array}
\]
The coefficient signs make all unmentioned entries in the sign vectors
immediate.  The intermediate nonautomatic entries use only the following
additional comparisons.  At \(w=10\), with \(y=10-3\pi\),
\[
 \tan y<\frac{2900}{4159}<\frac{170}{117}<10,
 \qquad
 \tan y>y>\frac47>\frac{890}{21789};
\]
substitution in \(P_1,P_3,P_5\) gives the signs \(+,-,-\), and
\(P_4=7P_3-w^2P_2<0\).  At \(w=103/10\), with
\(u=7\pi/2-w\),
\[
 \tan u>\frac{69}{100}>\frac{10}{103},
\]
which gives \(\mathsf j_1(w)>0\).  At \(w=23/2\), with
\(u=4\pi-w\),
\[
 \tan u>\frac{546377}{375000}>
 \frac{3153151489}{2276518664}>
 \frac{224020}{186301}>\frac{138}{517};
\]
besides the target \(\mathsf j_7\) sign, the last two thresholds give
\(\mathsf j_4(w)<0\) and \(\mathsf j_2(w)>0\), and the recurrence supplies
the intervening signs.  Finally, at \(w=129/10\), with \(y=w-4\pi\),
\[
 \tan y<\frac{1700}{4711}<
 \frac{177800829}{427700150}<
 \frac{651063}{327820}<\frac{129}{10};
\]
besides the target \(\mathsf j_5\) sign, this proves
\(\mathsf j_3(w)>0\) and \(\mathsf j_1(w)<0\), with the remaining signs
again immediate from the recurrence.  Thus every sign used in the count
induction is explicitly certified.

All comparisons are between positive rationals and are checked by cross
multiplication.  In every table row the target component is
\(Z_\ell(w)=n-1\), and the displayed nonzero sign proves
\(\mathsf j_\ell(w)\ne0\).  With \(j_{\nu,0}:=0\), it follows that
\[
 j_{\ell+1/2,n-1}<w<j_{\ell+1/2,n}.
\]
Thus the unique next zero has the asserted index, and the ledger proves
\begin{equation}
\begin{gathered}
 j_{3/2,2}>77/10,\quad j_{5/2,2}>9,\quad
 j_{7/2,2}>103/10,\quad j_{11/2,2}>129/10,\\
 j_{3/2,3}>21/2,\quad j_{5/2,3}>61/5,\\
 j_{13/2,1}>10,\quad j_{15/2,1}>23/2.
\end{gathered}
 \label{cm:eq-internal-zero-registry}
\end{equation}

For the remaining first zeros we use the two strict lower bounds displayed
as (i) and (ii) in the published abstract of Lorch~\cite{Lorch1993}, both
valid for the complete stated range \(-1<\nu<\infty\):
\begin{align}
 j_{\nu,1}^2&>(\nu+1)(\nu+5), \tag{L1}\label{cm:eq-lorch-L1}\\
 j_{\nu,1}^2&>
 \frac{24(\nu+1)^2}
  {1-2\nu+\sqrt{(2\nu+3)(2\nu+11)}}-2(\nu^2-1),
 \tag{L2}\label{cm:eq-lorch-L2}
\end{align}
At \(\nu=\ell+1/2\), the right side of \eqref{cm:eq-lorch-L2} is
\[
 R_\ell=\frac{3(2\ell+3)\sqrt{(\ell+2)(\ell+6)}
 -2\ell^2+\ell+6}{4}.
\]
Exact positive-side squaring gives
\begin{center}
\small
\begin{tabular}{c|c|c}
\(\ell\)&first-zero wall&exact strict witness\\ \hline
2&\(51/10\)&\((\mathrm{L1}):\ 105/4-(51/10)^2=6/25\)\\
3&\(13/2\)&\(81^2\cdot5-178^2=1121\)\\
4&\(15/2\)&\(66^2\cdot15-247^2=4331\)\\
5&\(87/10\)&\(975^2\cdot77-8544^2=198189\)\\
8&\(71/6\)&\(1026^2\cdot35-6067^2=35171\)\\
9&\(64/5\)&\(1575^2\cdot165-20059^2=6939644\)\\
10&\(69/5\)&\(3450^2\cdot3-5911^2=767579\).
\end{tabular}
\end{center}
Together with \eqref{cm:eq-internal-zero-registry}, this yields the complete
first-mode walls
\begin{equation}
\begin{array}{c|cccccccccc}
\ell&1&2&3&4&5&6&7&8&9&10\\ \hline
\beta_{\ell,1}&4&51/10&13/2&15/2&87/10&10&23/2&71/6&64/5&69/5.
\end{array}
 \label{cm:eq-first-zero-table}
\end{equation}
The \(\ell=1\) wall follows internally from \(\tan x=x\).  Its first
positive solution lies uniquely in \((\pi,3\pi/2)\), since, for
\(f(x)=\tan x-x\),
\[
 f(\pi)=-\pi<0,\qquad
 \lim_{x\uparrow3\pi/2}f(x)=+\infty,\qquad
 f'(x)=\tan^2x>0.
\]
Moreover \(\pi<4<3\pi/2\).  Put \(s=4-\pi\).  Then
\[
 0<s<4-\frac{333}{106}=\frac{91}{106}<1,
 \qquad
 \tan4=\tan s<\frac{s}{1-s^2/2}<2<4,
\]
where the penultimate inequality is equivalent to
\(2-s-s^2=(1-s)(s+2)>0\).  Hence the first root is strictly larger
than \(4\).
Applying \eqref{cm:eq-angular-propagation} to the higher-radial seeds gives,
for later use,
\begin{align}
 j_{p+1/2,2}^2&>9409/100+p(p+1) &&(p\geq3),\label{cm:eq-prop-second-a}\\
 j_{p+1/2,3}^2&>3571/25+p(p+1) &&(p\geq2),\label{cm:eq-prop-third}\\
 j_{p+1/2,2}^2&>13641/100+p(p+1) &&(p\geq5).
 \label{cm:eq-prop-second-b}
\end{align}
Equations \eqref{cm:eq-shell-ball}--\eqref{cm:eq-prop-second-b} transfer
every displayed wall to the shell at the same radial index.


\subsection{The two-sided staircase}

Put
\[
 \rho_0=\frac1{\sqrt{337}},\qquad d=\frac{\sqrt{337}}2,
 \qquad p_1(\rho)=\sqrt{4z_\rho^2+2}.
\]

\begin{lemma}[High side through \(k_6\)]
\label{cm:lem-k6}
For \(\rho_c\leq\rho\leq39/50\) and
\(z_\rho\leq K\leq k_6(\rho)\), one has \(N_D<W\).
\end{lemma}

\begin{proof}
For \(\rho\leq1/4\), the only possible new modes beyond the first modes
\(\ell=0,\ldots,4\) are \((0,2),(1,2)\); the first \(\ell=5\) mode is
delayed by \(j_{11/2,1}>17/2\).  For \(\rho\geq1/4\),
\(z>4\) makes every second mode absent through \(k_6\), and first modes
\(\ell=0,\ldots,5\) are the complete inventory.  The cumulative first-mode
caps are \(1,4,9,16,25,36\); the two possible second modes give caps
\(26,29\).  Their entry walls are
\[
\begin{array}{c|c}
(1,1)&k_1\\
(2,1)&\max\{k_2,51/10\}\\
(3,1)&\max\{k_3,13/2\}\\
(4,1)&\max\{k_4,15/2\}\\
(5,1)&\max\{k_5,17/2\}\\
(0,2)&2z\\
(1,2)&\max\{p_1,77/10\}.
\end{array}
\]
Here \(17/2<87/10\).  We deliberately weaken the registry wall
\(\beta_{5,1}=87/10\) to \(17/2\); this can only enlarge the possible-mode
inventory and it matches the fixed payment split below.
The initial cap is paid uniformly because
\[
 W(\rho,z_\rho)=\frac{2\pi^2}{9}
 \frac{1+\rho+\rho^2}{(1-\rho)^2}
 >\frac{2\pi^2}{9}>2>1.
\]
For a moving wall \(T^2=az^2+b\), \(W(\rho,T)\) increases whenever
\(a\pi^2(1+\rho)>b\rho^2(1-\rho)^2\), which holds here because
\(b/a\leq42\).  Fixed walls decrease with \(\rho\).  Exact payments are
obtained by the splits
\[
 (k_1,1/8,4),\ (k_2,3/10,9),\ (k_3,1/2,16),\
 (k_4,1/2,25),\ (k_5,13/25,36),\ (p_1,1/5,29).
\]
Writing
\(A=a(333/106)^2/(1-r)^2+b\) and
\(B=7(1-r^3)/99\), each moving payment is the positive rational check
\(B^2A^3-C^2>0\).  Here \(A<T(r)^2\) uses
\(\pi>333/106\), while
\[
 B=\frac7{99}(1-r^3)
 <\frac{2}{9\pi}(1-r^3)
\]
uses \(1/\pi>7/22\), equivalently \(\pi<22/7\).  In the displayed
order the numerators are
\[
\begin{gathered}
153864824754340538785,
1399810848073457853053,
137027505353736631,\\
2507119282010251109,
92672404686738742434741,
373255772037478993002833,
\end{gathered}
\]
over positive denominators.  The corresponding fixed-side margins are
positive as well; for example they are
\(1387329/11000000,6277/6336,775/704,452843/343750\), and
\(849901/281250\) at the nontrivial splits.  Finally
\(W(\rho_c,2z)>26\).  Every possible mode is therefore paid, including
coincident entries, and every equality remains assigned to the lower cap.
 Parts~I, II, and IV of the \emph{Purely Analytic Supplement} map every
 printed moving numerator and fixed margin to its wall and cap.
\end{proof}

\begin{lemma}[Lower side through \(d\)]
\label{cm:lem-lower-d}
For \(\rho_0<\rho<\rho_c\) and
\((2\rho)^{-1}<K\leq d\), one has \(N_D<W\).
\end{lemma}

\begin{proof}
Because \(3z>d\), every \(n\geq3\) mode is absent.  The zero registry and
angular propagation exclude first modes \(\ell\geq6\) and second modes
\(\ell\geq3\).  The complete cap table is
\[
\begin{array}{c|c}
\text{open-lower, closed-upper range}&N_D\text{ cap}\\ \hline
L<K\leq4&1\\
4<K\leq51/10&4\\
51/10<K\leq13/2,\ K\leq2z&9\\
51/10<K\leq13/2,\ K>2z&10\\
13/2<K\leq15/2,\ K\leq2z&16\\
13/2<K\leq15/2,\ K>2z&17\\
15/2<K\leq77/10&26\\
77/10<K\leq17/2&29\\
17/2<K\leq9&40\\
9<K\leq d&45.
\end{array}
\]
In fact \(13/2<2z<15/2\), so the cap-10 row is empty.  The fixed-wall
lower bounds, using \(\rho<\rho_c<7/50\), exceed their caps by
\[
 \frac{793292}{1546875},\frac{486236661}{1375000000},
 \frac{333100003}{99000000},\frac{329797}{88000},
 \frac{3590476297}{1125000000},\frac{327078887}{99000000},
 \frac{8805519}{1375000},
\]
respectively.  The initial cap is paid by
\(W(\rho_c,L(\rho_c))>19/6>1\), and the moving branch by
\(W(1/19,2z)>508/27>17\).  This proves every row, including \(K=d\).
 Parts~I and IV of the \emph{Purely Analytic Supplement} record each margin
 with its corresponding row and prove the moving \(2z\) payment.
\end{proof}

The upper-floor comparisons are strict.  On the lower side,
\(H_0(\rho_*)=(2\rho_*)^{-1}>d\), while
\(K_0>C_0/\omega_0>12>d\).  On the high side, \(K_0>12>k_6\) for
\(\rho\leq1/2\), while \(K_0>384>k_6\) for
\(1/2<\rho<39/50\).  Thus exact subtraction gives
\begin{equation}
\boxed{\begin{aligned}
\mathcal D_{19}={}&
\{\rho_*<\rho\leq\rho_0:(2\rho)^{-1}<K<U(\rho)\}\\
&\cup\{\rho_0<\rho<\rho_c:d<K<U(\rho)\}\\
&\cup\{\rho_c\leq\rho<39/50:k_6(\rho)<K<U(\rho)\}.
\end{aligned}}
\label{cm:eq-D19}
\end{equation}
At \(\rho=\rho_0\), \((2\rho_0)^{-1}=d\), so the would-be new lower
fiber is empty and the first component owns the equality slice.

\subsection{Combined closure of the lower components and the high staircase}

Set
\[
 \sigma:=\frac{3\omega_0}{4}.
\]
We make every comparison in the required ordering explicit.  The
alternating Machin sums give
\[
\begin{aligned}
 \pi
 &<16\left(\frac15-\frac1{3\cdot5^3}
 +\frac1{5\cdot5^5}-\frac1{7\cdot5^7}
 +\frac1{9\cdot5^9}\right)\\
 &\quad-4\left(\frac1{239}-\frac1{3\cdot239^3}\right)
 <\frac{355}{113},
\end{aligned}
\]
and the last gap is exactly
\[
 \frac{45167474711}{189820334689453125}>0.
\]
Since
\[
 3\cdot1101^2\cdot113^2-607^2\cdot355^2=1998482>0,
\]
we have
\[
 \sqrt3>\frac{607\cdot355}{1101\cdot113}
 >\frac{607\pi}{1101},
\]
and therefore
\begin{equation}
 \omega_0=\frac{\sqrt3}{2\pi}-\frac16>\frac{40}{367}.
 \label{cm:eq-omega-lower}
\end{equation}
On the other side,
\[
 25\cdot333^2-243\cdot106^2=41877>0
\]
and \(\pi>333/106\) imply \(9\sqrt3<5\pi\), hence
\(\omega_0<1/9\).  Moreover
\(2\cdot32^2-45^2=23>0\), so
\(\sqrt2>45/32\), \(C_0>7/4\), and
\[
 \rho_*=\frac{\omega_0}{2C_0}<\frac2{63}.
\]
Now all remaining comparisons are witnessed by
\[
\begin{gathered}
 4\cdot337<63^2,\qquad 18^2<337,\qquad
 \frac1{18}<\frac3{40}<\frac{3\omega_0}{4}
 <\frac1{12}<\frac7{51},\\
 \frac7{51}<\frac1{1+2\pi}<\frac17,\qquad
 \frac17<\frac{39}{50}.
\end{gathered}
\]
Here the first inequality on the second line uses \(\pi<22/7\), and the
second uses \(\pi>3\).  Thus, without decimal comparison,
\[
 \rho_*<\rho_0<\sigma<\rho_c<\frac{39}{50}.
\]

\begin{lemma}[Closure of both lower components of \(\mathcal D_{19}\)]
\label{cm:lem-lower-closure}
The strict inequality \(N_D<W\) holds on the first two components of
\eqref{cm:eq-D19}.
\end{lemma}

\begin{proof}
We divide the proof into a wedge, a finite staircase, and one remaining
floor cell.

For a low-interface tail beginning at \(x_r=r+1/2<\mu=\rho K\), put
\[
 n_r=\lfloor\mu-x_r\rfloor,\qquad
 q_r=x_r+n_r,\qquad p=\lfloor\omega_0K\rfloor.
\]
The concave-head/convex-tail split in
\eqref{eq:sc-tail-compensation}, now with its coarse first-shelf payment,
gives
\begin{equation}
 \frac{\mathcal T_r}{2}
 <\int_{x_r}^K A_{\rho,K}(z)\,dz
 +\delta_r-\frac{p-n_r}{4},
 \qquad
 0\leq\delta_r<
 \frac{2(\mu-q_r)^{5/2}}{15\sqrt\mu}.
 \label{cm:eq-lower-wedge-tail}
\end{equation}
When \(n_r=0\), the concave head is absent and the formula follows from
the convex tail alone.  From \eqref{cm:eq-omega-lower} and
\(400\cdot337-367^2=111>0\),
\[
 \omega_0d>
 \frac{40}{367}\frac{\sqrt{337}}2
 =\frac{20\sqrt{337}}{367}>1.
\]
If
\(0<\rho\leq\sigma\), \(K>d\), and \(\rho K>1/2\), then with
\(x=\omega_0K>1\) and \(m=\lfloor x\rfloor\),
\[
 n_r\leq\left\lfloor\frac{3x}{4}-\frac12\right\rfloor<m=p.
\]
Since \(\delta_r<2\sqrt2/15<1/4\), every tail in
\eqref{cm:eq-lower-wedge-tail} is strict.  The hypotheses hold throughout
the first component of \(\mathcal D_{19}\) and on
\(\rho_0\leq\rho\leq\sigma\), \(K>d\).  This includes the complete
\(\rho=\rho_0\) fiber.

For \(\rho_0<\rho<\rho_c\) and \(K\leq\sqrt{114}\), the complete zero
registry and \eqref{cm:eq-angular-propagation} give the following exhaustive
inventory:
\begin{center}
\small
\begin{tabular}{c|l|c}
range (lower face open, upper face closed)&newly possible modes&cap\\ \hline
\(d<K\leq10,\ K\leq3z\)&first \(0{:}5\), second \(0{:}2\)&45\\
\(d<K\leq10,\ K>3z\)&preceding plus \((0,3)\)&46\\
\(10<K\leq81/8\)&plus \((6,1)\), possibly \((0,3)\)&59\\
\(81/8<K\leq21/2\)&plus \((3,2)\)&66\\
\(21/2<K\leq\sqrt{7073}/8\)&plus \((1,3)\)&69\\
\(\sqrt{7073}/8<K\leq\sqrt{114}\)&plus \((4,2)\)&78.
\end{tabular}
\end{center}
First modes \(\ell\geq7\), second modes \(\ell\geq5\), third modes
\(\ell\geq2\), and all \(n\geq4\) are absent.  In particular,
\[
 (81/8)^2+8=7073/64
\]
is the propagated \((4,2)\) wall.  Since \(\rho<\rho_c<1/7\), the
uniform estimate
\[
 W(\rho,K)>
 \frac{2}{9(22/7)}\left(1-\frac1{7^3}\right)K^3
\]
pays the successive caps with exact positive reserves
\[
 \frac{2908}{539},\quad\frac{6199}{539},\quad
 \frac{990435}{137984},\quad\frac{555}{44},\quad\frac{159}{44}.
\]
The first number pays even cap \(46\) from \(K>9\); the last pays cap
\(78\) from \(K>21/2\).  Hence \(K=\sqrt{114}\) is included.

It remains to consider \(\sigma<\rho<\rho_c\) and
\(K>\sqrt{114}\).  Put
\[
 p=\lfloor\omega_0K\rfloor,\qquad
 m=\lfloor\rho K-1/2\rfloor.
\]
Then \(p\geq1\).  Except on
\begin{equation}
 \mathcal B=\left\{\sigma<\rho<\rho_c,\quad K>\sqrt{114},\quad
 \rho K\geq5/2,\quad K<2/\omega_0\right\},
 \label{cm:eq-exceptional-B}
\end{equation}
one has \(m\leq2p-1\).  Indeed this is immediate when \(\rho K<5/2\);
when \(K\geq2/\omega_0\), one has \(p\geq2\).  Moreover,
\[
 \rho_c<\frac17<\frac{60}{367}<\frac{3\omega_0}{2},
\]
where the first and last comparisons were proved above and
\(367<420\).  Hence \(\rho K-1/2<2p\).

Away from a half-integer \(\mu\), all interface remainders in
\eqref{cm:eq-lower-wedge-tail} are equal.  Summing over
\(r=0,\ldots,m\) gives the average floor reserve
\[
 \frac14\left(p-\frac m2\right)\geq\frac18.
\]
For \(m=0\), use \(\delta<1/4\); for \(m\geq1\), \(\mu>3/2\) and
\(\delta<2/(15\sqrt{3/2})<1/8\).  At
\(\mu\in\mathbb Z+1/2\), assign the interface to the high side; then the
remainder is zero and strictness remains in the nonzero convex tail.
Consequently the sum of all tail excesses is strictly negative.  This is
exactly the aggregate hypothesis \eqref{eq:sc-aggregate-tail-target} of
Lemma~\ref{lem:sc-weighted-scaffold}, which closes the nonexceptional
region; no per-tail inequality is inferred from the averaged reserve.

On \(\mathcal B\), \(\omega_0>40/367\), so \(K<367/20=:R\).  The
turning estimate
\[
 A_{\rho,K}(z)\leq G_K(z)<\frac{(K-z)^{3/2}}{3\sqrt K}
\]
and monotonicity in \(K\) reduce the proxy count at
\(z_\ell=\ell+1/2\) to
\[
 \frac{(357-20\ell)^3}{1321200}
 <\left(c_\ell+\frac34\right)^2.
\]
The integer numerators of these strict comparisons are
\begin{center}
\small
\begin{tabular}{c|rrrrrrrrrrrrrrr}
\(\ell\)&0&1&2&3&4&5&6&7&8&9&10&11&12&13&14\\ \hline
\(c_\ell\)&6&5&5&4&4&3&3&3&2&2&1&1&1&1&0
\end{tabular}
\end{center}
and, in the same order,
\[
\begin{gathered}
14697882,5409422,11827162,3611502,8555642,
1604782,5267322,8361062,\\
2346202,4446342,176282,1474822,2444562,3133502,286642.
\end{gathered}
\]
The \(\ell=14\) row and monotonicity remove the entire angular tail.
Therefore the full multiplicity cap is
\[
 \sum_{\ell=0}^{13}(2\ell+1)c_\ell=395.
\]
But \(\rho K\geq5/2\), \(\rho<\rho_c<11/80\), and
\(1/\pi>113/355\) give
\[
 W(\rho,K)>
 \frac{160293325}{378004}
 =395+\frac{10981745}{378004}>395.
\]
Thus \(\mathcal B\) is closed as well.  The face \(\rho K=5/2\) is
\(\mathcal B\)-owned and the complementary branch owns
\(K=2/\omega_0\).  All lower
components of \(\mathcal D_{19}\) are proved.
\end{proof}

We now turn to the high component.  Define the strict channel proxy
\begin{equation}
 b_{\ell,n}(z)^2=
 \max\{n^2z^2+\ell(\ell+1),\,\beta_{\ell,n}^2\},
 \label{cm:eq-channel-proxy}
\end{equation}
For the radial channel we set, for every \(n\geq1\),
\[
 \beta_{0,n}:=j_{1/2,n}=n\pi,
\]
the last identity following from
\(J_{1/2}(x)=\sqrt{2/(\pi x)}\sin x\).  For \(\ell\geq1\), the
first-mode \(\beta\)'s are in
\eqref{cm:eq-first-zero-table} and
\[
\begin{aligned}
 \beta_{1,2}&=77/10,&
 \beta_{2,2}&=9,&
 \beta_{3,2}&=103/10,&
 \beta_{4,2}&=\sqrt{11409}/10,\\
 \beta_{5,2}&=129/10,&
 \beta_{1,3}&=21/2,&
 \beta_{2,3}&=61/5.
\end{aligned}
\]
By \eqref{cm:eq-product-bound} and \eqref{cm:eq-shell-ball}, a mode can
be counted only if \(b_{\ell,n}(z)<K\).
For \(\ell\geq1\), every displayed registry constant is a strict lower
wall,
\begin{equation}
 \beta_{\ell,n}<j_{\ell+1/2,n}\leq k_{\ell,n}(\rho),
 \label{cm:eq-proxy-versus-zero}
\end{equation}
whereas \(\beta_{0,n}=j_{1/2,n}=n\pi\) is exact.  Thus equality at a
\(\beta\)-wall is proxy equality, not a claim that an actual eigenfrequency
equals the rational wall; the proxy deliberately permits a mode no later
than the spectrum does.

\begin{lemma}[Closed high staircase through \(k_{11}\)]
\label{cm:lem-k11}
For \(\rho_c\leq\rho\leq39/50\) and
\(z_\rho\leq K\leq k_{11}(\rho)\), one has \(N_D<W\).
\end{lemma}

\begin{proof}
Lemma~\ref{cm:lem-k6} already proves the claim for
\(z_\rho\leq K\leq k_6(\rho)\), including the common face.  It therefore
suffices to treat \(k_6(\rho)<K\leq k_{11}(\rho)\).
The exhaustive checkpoint inventories are
\begin{center}
\small
\begin{tabular}{c|l}
checkpoint&possible modes\\ \hline
\(k_7\)&first \(0{:}6\), second \(0{:}1\), no \(n\geq3\)\\
\(k_8\)&first \(0{:}7\), second \(0{:}3\), no \(n\geq3\)\\
\(k_9\)&first \(0{:}8\), second \(0{:}4\), no \(n\geq3\)\\
\(k_{10}\)&first \(0{:}9\), second \(0{:}5\), third \(0{:}1\), no \(n\geq4\)\\
\(k_{11}\)&first \(0{:}10\), second \(0{:}4\), third \(0{:}1\), no \(n\geq4\).
\end{tabular}
\end{center}
Here is the promised exhaustion check.  First channels with \(\ell\geq m\)
are excluded at \(k_m\) by their product wall.  For the first omitted second
channel, the product and ball walls would force the incompatible strict
bounds
\[
\begin{array}{c|ccccc}
m&7&8&9&10&11\\ \hline
\text{first omitted }\ell&2&4&5&6&5\\
\text{product wall}&x<50/3&x<52/3&x<20&x<68/3&x<34\\
\text{ball wall}&x>25&x>4209/100&x>7641/100&
x>6841/100&x>3441/100.
\end{array}
\]
The \(m=10\) ball wall is the angular propagation
\eqref{cm:eq-prop-second-b}; the other four are displayed registry walls.
For larger \(\ell\), angular propagation raises the ball lower bound while
the product upper bound decreases.

For \(n=3\), the product wall gives
\(x\leq m(m+1)/8\leq45/4<(193/53)^2<x\) when \(m\leq9\).
At \(m=10\), every \(\ell\geq2\) gives
\(x\leq(110-6)/8=13<(193/53)^2\).  At \(m=11\), the first omitted
third channel \(\ell=2\) would require
\[
 x<\frac{132-6}{8}=\frac{63}{4},
 \qquad
 x>\left(\frac{61}{5}\right)^2-132=\frac{421}{25},
\]
and higher angular indices are again excluded by propagation.  Finally,
for \(n\geq4\), already \(m=11,\ell=0,n=4\) would require
\(x<132/15=44/5<(193/53)^2\), and all higher indices are stricter.
This proves the five inventories without an unprinted finite search.

Writing \(H\) for the highest first channel and \(h\) for the highest
second channel, the block cap is \((H+1)^2+(h+1)^2\), with an additional
\(1\) or \(4\) for third modes.  The complete cap tables at \(k_8\) and
\(k_9\) are
\[
\begin{array}{c|ccccc}
k_8:H\backslash h&\varnothing&0&1&2&3\\ \hline
7&64&65&68&73&80\\
6&49&50&53&-&-\\
5&36&37&40&45&52\\
\leq4&25&26&29&34&41
\end{array}
\qquad
\begin{array}{c|cccccc}
k_9:H\backslash h&\varnothing&0&1&2&3&4\\ \hline
8&81&-&-&-&-&-\\
7&64&65&68&73&-&-\\
6&49&50&53&58&65&74\\
5&36&37&40&45&52&61\\
\leq4&25&26&29&34&41&50.
\end{array}
\]
The dashes follow from incompatible necessary inequalities, not from a
numerical sample.  At \(k_{10}\) the cap lists are
\[
 100;\quad81,97;\quad64,80,89,100;\quad74,69,78,85,89,
\]
and at \(k_{11}\) they are
\[
 121;\quad100,125;\quad106,110;\quad93.
\]

We now give the missing order-free exhaustion of all events.  Let
\(\mathcal I\) be the finite union of the five checkpoint inventories above
and define both the strict and closed post-event proxies
\[
 C_z(K)=\sum_{(\ell,n)\in\mathcal I}(2\ell+1)
 \mathbf1_{\{b_{\ell,n}(z)<K\}},\qquad
 \widehat C_z(K)=\sum_{(\ell,n)\in\mathcal I}(2\ell+1)
 \mathbf1_{\{b_{\ell,n}(z)\leq K\}}.
\]
Thus \(C_z\leq\widehat C_z\).  Put
\[
 \mathcal B_m=\{k_{m-1}\leq K<k_m\}\quad(7\leq m\leq11).
\]
After all modes on a combined event wall have been inserted, let \(H\) and
\(h\) be the largest occupied first- and second-mode angular indices, with
\(h=-1\) when no second mode is present, and let \(T\in\{0,1\}\) record
whether any third mode is present.  Contiguity is not assumed: the odd-sum
identity always gives the safe full-block cap
\begin{equation}
 \widehat C_z(K)\leq (H+1)^2+(h+1)^2+4T.
 \label{cm:eq-closed-post-event-cap}
\end{equation}

For \(j\leq10\), set \(F_j(\rho)=W(\rho,k_j(\rho))\), and define
\[
 \mathsf P(C;c,r)=\frac7{99}(1-r^3)c^3-C,
 \qquad
 \mathsf Q_j(c,r)=
 \left(\frac{333}{106(1-r)}\right)^2+j(j+1)-c^2.
\]
There are three payment rules.
\begin{enumerate}
\item \(\mathsf G_j(C)\): the increasing checkpoint payment satisfies
      \(F_j>C\).
\item \(\mathsf L(C;c,r)\): the state forces \(K\geq c\) and \(\rho<r\),
      while \(\mathsf P(C;c,r)>0\).
\item \(\mathsf B_j(C;c,r)\): the state forces \(K\geq c\), while
      \(\mathsf P(C;c,r)>0\) and \(\mathsf Q_j(c,r)>0\).
\end{enumerate}
The localized rule follows from
\[
 W(\rho,K)>
 \frac7{99}(1-\rho^3)K^3>
 \frac7{99}(1-r^3)c^3>C.
\]
For the bridge rule, this argument applies when \(\rho\leq r\).  When
\(\rho\geq r\), the band lower wall gives \(K\geq k_j(\rho)\);
\(\mathsf Q_j>0\) gives \(k_j(r)>c\); and \(F_j\) is strictly increasing
because
\[
 \pi^2(1+\rho)-j(j+1)\rho^2(1-\rho)^2
 >9-\frac{110}{16}>0.
\]
Moreover
\[
 F_j(r)=W(r,k_j(r))>W(r,c)>
 \frac7{99}(1-r^3)c^3>C.
\]
Hence \(W(\rho,K)\geq F_j(\rho)\geq F_j(r)>C\).  The same proof owns the
split face \(\rho=r\).

For completeness, the four global rules used below follow at
\(\rho_c\) from \(z_{\rho_c}=\pi+1/2>193/53\) and
\[
 W(\rho_c,K)>\frac{38}{539}K^3.
\]
The two last columns are respectively the positive square gap proving
\(k_j(\rho_c)>t\) and the resulting payment reserve:
\begin{center}
\scriptsize
\begin{tabular}{c|c|c|c|c}
\(j\)&\(C\)&\(t\)&\((193/53)^2+j(j+1)-t^2\)&
\((38/539)t^3-C\)\\ \hline
7&40&\(208/25\)&\(67049/1755625\)&\(5083656/8421875\)\\
8&53&\(923/100\)&\(1901439/28090000\)&
\(656778873/269500000\)\\
9&73&\(254/25\)&\(61431/1755625\)&\(7911557/8421875\)\\
10&89&\(111/10\)&\(14211/280900\)&\(1999489/269500\)
\end{tabular}
\end{center}

The following table is the complete state-to-payment map.  In each band,
\emph{remaining} means the literal complement of the other displayed
states; the cap in that row is an upper bound for every such state.  Thus
the rows are disjoint and exhaustive, not a list of representative cases.
\begingroup
\scriptsize
\renewcommand{\arraystretch}{1.12}
\begin{longtable}{@{}c p{5.8cm} c p{5.6cm}@{}}
\toprule
band&post-event state&cap&payment certificate\\
\midrule
\endfirsthead
\toprule
band&post-event state&cap&payment certificate\\
\midrule
\endhead
\(\mathcal B_7\)&remaining&26&\(\mathsf B_6(26;73/10,1/8)\)\\
&\(H=4,\ h=1\)&29&\(\mathsf L(29;77/10,3/10)\)\\
&\(H=5,\ h=-1\)&36&\(\mathsf B_6(36;87/10,1/2)\)\\
&\(H=5,\ h\geq0\)&40&\(\mathsf L(40;87/10,3/10)\)\\
&\(H=6,\ h=-1\)&49&\(\mathsf B_6(49;10,3/5)\)\\
&\(H=6,\ h\geq0\)&53&\(\mathsf L(53;10,3/10)\)\\
\midrule
\(\mathcal B_8\)&remaining&40&\(\mathsf G_7(40)\)\\
&\(H=5,\ h=2\)&45&\(\mathsf L(45;9,1/3)\)\\
&\(H\in\{4,5\},\ h=3\)&52&\(\mathsf L(52;103/10,3/10)\)\\
&\(H=6,\ h=-1\)&49&\(\mathsf B_7(49;10,3/5)\)\\
&\(H=6,\ h\geq0\)&53&\(\mathsf L(53;10,3/8)\)\\
&\(H=7,\ h=-1\)&64&\(\mathsf B_7(64;23/2,2/3)\)\\
&\(H=7,\ h\geq0\)&80&\(\mathsf L(80;23/2,3/8)\)\\
\midrule
\(\mathcal B_9\)&remaining&53&\(\mathsf G_8(53)\)\\
&\(H=6,\ h=2\)&58&\(\mathsf L(58;10,5/12)\)\\
&\(H=6,\ h=3\)&65&\(\mathsf L(65;103/10,2/5)\)\\
&\(H\in\{5,6\},\ h=4\)&74&\(\mathsf L(74;207/20,7/20)\)\\
&\(H=7,\ h=-1\)&64&\(\mathsf B_8(64;23/2,2/3)\)\\
&\(H=7,\ h\geq0\)&73&\(\mathsf L(73;23/2,3/7)\)\\
&\(H=8\)&81&\(\mathsf B_8(81;71/6,2/3)\)\\
\midrule
\(\mathcal B_{10}\)&remaining&73&\(\mathsf G_9(73)\)\\
&\(H=6,\ h=4,\ T=0\)&74&\(\mathsf L(74;53/5,3/7)\)\\
&\(H=6,\ h=4,\ T=1\)&78&\(\mathsf L(78;53/5,4/25)\)\\
&\(H=6,\ h=5,\ T=0\)&85&\(\mathsf L(85;11,2/5)\)\\
&\(H\in\{5,6\},\ h=5,\ T=1\)&89&\(\mathsf L(89;11,4/25)\)\\
&\(H=7,\ h=3\)&80&\(\mathsf L(80;23/2,1/2)\)\\
&\(H=7,\ h=4\)&89&\(\mathsf L(89;23/2,3/7)\)\\
&\(H=7,\ h=5\)&100&\(\mathsf L(100;23/2,2/5)\)\\
&\(H=8,\ h=-1\)&81&\(\mathsf B_9(81;71/6,3/5)\)\\
&\(H=8,\ h\geq0\)&97&\(\mathsf L(97;71/6,1/2)\)\\
&\(H=9\)&100&\(\mathsf B_9(100;64/5,16/25)\)\\
\midrule
\(\mathcal B_{11}\)&remaining&89&\(\mathsf G_{10}(89)\)\\
&\(H=7,\ h=4,\ T=1\)&93&\(\mathsf L(93;23/2,1/4)\)\\
&\(H=8,\ h\geq2,\ T=0\)&106&\(\mathsf B_{10}(106;71/6,3/7)\)\\
&\(H=8,\ h\geq2,\ T=1\)&110&\(\mathsf L(110;71/6,1/4)\)\\
&\(H=9,\ h=-1\)&100&\(\mathsf B_{10}(100;64/5,3/5)\)\\
&\(H=9,\ h\geq0\)&125&\(\mathsf L(125;64/5,8/15)\)\\
&\(H=10\)&121&\(\mathsf B_{10}(121;69/5,13/20)\)\\
\bottomrule
\end{longtable}
\endgroup

The four global rows are the exact checkpoint payments printed above.
For every localized row, occupancy of the named block gives \(K\geq c\);
its product inequality below the upper checkpoint gives the stated
\(\rho<r\).  For example,
\[
 z_{3/10}^2-\frac{56}{3}>
 \frac{608779}{412923}>0,\qquad
 z_{7/20}^2-\frac{70}{3}>
 \frac{36230}{1424163}>0,\qquad
 z_{8/15}^2-44>
 \frac{725209}{550564}>0.
\]
The remaining eleven localization differences, all values of
\(\mathsf P\), and all bridge values of \(\mathsf Q_j\) are printed as
positive fractions in Part~III of the supplement.  The two tight bridge
checks are still strictly positive:
\[
 \mathsf Q_{10}(71/6,3/7)=\frac{317585}{1617984},
 \qquad
 \mathsf Q_{10}(69/5,13/20)=\frac{426449}{3441025}.
\]
The first row uses \(\rho\geq\rho_c>1/8\), so only the increasing-checkpoint
side of its bridge is needed.

This proves \(W>\widehat C_z\) throughout every half-open band.  Arbitrary crossings
of fixed and moving event walls do not require a separate ordering table:
after each event the resulting triple \((H,h,T)\) lies in exactly one row.
At a coincidence, \(\widehat C_z\) inserts every entering multiplicity
before selecting that row, so the combined jump is paid.  At the wall
itself the original strict proxy has only the smaller pre-event cap.  A
checkpoint event is assigned post-event to the following half-open band.
At the final face \(K=k_{11}\), the original strict proxy omits the
\((11,1)\) equality channel; the same \(m=11\) state restrictions and
payments therefore give \(C_z(k_{11})<W(\rho,k_{11})\) directly.  Hence every
event, equality face, and intervening constant-count cell is exhausted.
The checkpoint inventories and \eqref{cm:eq-channel-proxy} therefore give
\[
 N_D(A_{\rho,1},K^2)\leq C_z(K)\leq\widehat C_z(K)<W(\rho,K)
\]
for \(k_6\leq K<k_{11}\), and
\[
 N_D(A_{\rho,1},k_{11}^2)\leq C_z(k_{11})<W(\rho,k_{11})
\]
on its final face.

There is one conservative row requiring explicit treatment.  Angular
propagation gives
\[
 j_{9/2,2}^2>\left(\frac{103}{10}\right)^2+8
 =\frac{11409}{100}.
\]
If \((4,2)\) were counted at \(k_9\), this ball wall would force
\(z^2>2409/100\), whereas the product wall forces \(z^2<70/3\); since
\(2409/100-70/3=227/300>0\), the printed cap-74 cell is actually empty.
If one deliberately retains the weaker conservative cell, counting the
mode forces \(K>207/20\) and \(\rho<7/20\).  Therefore
\begin{equation}
 W(\rho,K)>
 \frac{52823261673}{704000000}
 =74+\frac{727261673}{704000000}>74.
 \label{cm:eq-cap74}
\end{equation}
The lower bound in \eqref{cm:eq-cap74} uses
\(1/\pi>7/22\), i.e. \(2/(9\pi)>7/99\).
This is the corrected cap-74 payment; it is smaller than the generic
\(k_9\) summary reserve and must not be hidden by a claim that all omitted
cells have larger reserves.

At \(K=k_m\), \(6\leq m\leq10\), the strict proxy excludes the
\((m,1)\) product-equality mode, while the following band's closed cap
already pays it.  At \(K=k_{11}\), that equality channel remains absent
from the original strict proxy.  At a fixed
\(\beta\)-wall, \eqref{cm:eq-proxy-versus-zero} shows that the actual
spectral mode has not yet entered, while the closed cap again includes it.
Thus the order-free domination covers both kinds of equality, as well as
their coincidences, and proves the lemma.
\end{proof}

The finite analytic support---covering the zero/index registry, the lower
staircases, the closed post-event theorem, optical constants, and exact
subtraction---uses no floating-point premise.  Every sign decision is a
rational cross multiplication or a positive-side squaring.  Parts~I--VI of
the \emph{Purely Analytic Supplement} print the hand witnesses.  Their
historical 611-row census is retained for traceability; the final
38-row post-event state theorem is displayed separately in Part~III and is
the argument used above.

\fi

\section{Global two-ratio closure}
\label{sec:global-two-owner}

\subsection{Ratio-sharp defect theorem below the optical strip}

For $a>0$, let
\[
 G_a(z):=\frac1\pi\left(\sqrt{a^2-z^2}
 -z\arccos\frac za\right)\quad(0\leq z<a),
 \qquad G_a(z):=0\quad(z\geq a),
\]
and put
\[
 A(z):=G_K(z)-G_{\rho K}(z),
 \qquad T:=\frac{(1-\rho)K}{\pi}.
\]
The strict phase reduction and the inverse-action layer cake give
\begin{equation}
 N_D(A_{\rho,1},K^2)\leq P,\qquad
 P=\sum_{n-1/4<T}M(R(n-1/4))^2,\qquad
 W=\int_0^TR(t)^2\,dt,
 \label{rs-summary:eq:layer-cake}
\end{equation}
where $R$ is the decreasing inverse of $A$ and $M(r)$ counts
half-integers strictly below $r$.  Thus
\[
 W-P=\mathcal D_{\rm rad}-\mathcal E_{\rm ang}.
\]

Write $\theta=\arccos\rho$ and
$N=\#\{n\geq1:n-1/4<T\}$.  The exact action derivative satisfies
$-A'\leq\theta/\pi$.  One disjoint left block per active layer therefore
gives the ratio-sharp estimate
\begin{equation}
 \mathcal E_{\rm ang}
 <\frac{(1-\rho^2)K^2}{6}
 -N\left(\frac{3\pi}{8\theta}-\frac14\right).
 \label{rs-summary:eq:angular}
\end{equation}
Strictness includes a half-integer angular wall.

For $0<\rho\leq39/50$ and $K\geq\pi/(1-\rho)$, the curvature interface
$\tau=KG_1(\rho)$ satisfies $\tau>1/4$.  Replacing the shifted-sawtooth
primitive by its $C^1$ tangent envelope
\[
 L_\sharp(t)=
 \begin{cases}t^2/2,&t\leq1/4,\\t/4-1/32,&t\geq1/4,\end{cases}
\]
and retaining the decreasing-branch Stieltjes remainder gives
\begin{equation}
 \mathcal D_{\rm rad}\geq
 B_0(K)-\frac{\pi\rho K}{4(1-\rho)}
 +\frac{\pi\rho K}{4\arccos\rho},
 \qquad
 B_0(K):=\int_0^{1/4}R_B(t)^2\,dt,
 \label{rs-summary:eq:radial}
\end{equation}
where $R_B$ is the inverse ball action.  The turning estimate
$G_1(s)\geq2\sqrt2(1-s)^{3/2}/(3\pi)$ yields
\begin{equation}
 B_0(K)\geq\frac{K^2}{4}-\alpha K^{4/3}+\beta K^{2/3},
 \quad
 \alpha=\frac{6C}{5\,4^{5/3}},\quad
 \beta=\frac{3C^2}{7\,4^{7/3}},\quad
 C=\frac{(3\pi)^{2/3}}2.
 \label{rs-summary:eq:ball}
\end{equation}

Since $N\geq1$ and
$\arccos\rho\leq(\pi/2)\sqrt{1-\rho}$, it remains to prove positivity of
\begin{equation}
 \begin{aligned}
 \underline{\mathcal G}(\rho,K):={}&
 \frac{1+2\rho^2}{12}K^2-\alpha K^{4/3}+\beta K^{2/3}
 -\frac{\pi\rho K}{4(1-\rho)}\\
 &+\frac{\rho K}{2\sqrt{1-\rho}}
 +\frac3{4\sqrt{1-\rho}}-\frac14.
 \end{aligned}
 \label{rs-summary:eq:scalar}
\end{equation}
On the base $K=\pi/(1-\rho)$, the substitution
$x=(1-\rho)^{-1/6}$ reduces \eqref{rs-summary:eq:scalar} to one explicit
degree-nine expression.  The detailed support module proves it positive
from $157/50<\pi<22/7$, positive-side cubing, and
\[
 F(1)>\frac14,\qquad F'(1)>-\frac52,\qquad F''(x)>14.
\]
It also proves that the base $K$-derivative is positive and
\[
 \partial_K^2\underline{\mathcal G}
 >\frac{47447}{443682}>0.
\]
These are finite exact inequalities, not sampled or interval premises.

\begin{theorem}[Global low- and middle-ratio theorem]
\label{rs:thm:global-proxy}
For
\[
 0<\rho\leq\frac{39}{50},
 \qquad K\geq\frac\pi{1-\rho},
\]
one has
\[
 N_D(A_{\rho,1},K^2)\leq P<W(\rho,K).
\]
\end{theorem}

\begin{proof}
Equations \eqref{rs-summary:eq:angular}--\eqref{rs-summary:eq:ball}
give
\[
 W-P>\underline{\mathcal G}(\rho,K).
\]
The exact scalar closure just described proves the right side positive on
the stated domain.  All intermediate identities and the rational
convexity proof are given in the accompanying analytic supplement.
\end{proof}


\subsection{All-frequency optical theorem}

\begin{lemma}[All-frequency optical theorem]
\label{cm:lem-optical}
For \(39/50\leq\rho<1\) and \(K\geq0\),
\(N_D(A_{\rho,1},K^2)\leq W(\rho,K)\), with strict inequality for
\(K>0\).
\end{lemma}

\begin{proof}
Put \(\varepsilon=1-\rho\leq11/50\),
\(a=\varepsilon K\), and \(c=1126/625\).  We use two inclusive screens
meeting at \(a=c/\varepsilon\).

For the product screen, the strict min--max condition
\eqref{thin:eq-product-condition} implies that a counted mode satisfies
\[
 n^2\pi^2+\varepsilon^2\ell(\ell+1)<a^2.
\]
There is no mode for \(a\leq\pi\).  For \(a>\pi\), write
\(a/\pi=N+t\), with \(N\geq1\), \(0<t\leq1\), assigning
\((N,t)=(m-1,1)\) at \(a=m\pi\).  If
\[
 b_n=\sqrt{a^2-n^2\pi^2},\qquad
 M_n=\left\lceil\sqrt{(b_n/\varepsilon)^2+1/4}-1/2\right\rceil,
\]
the exact angular multiplicity is bounded by
\[
 \mathcal P_\varepsilon(a)=\sum_{n=1}^N M_n^2.
\]
Set
\[
 I_0(a)=\int_0^{a/\pi}(a^2-\pi^2x^2)\,dx
       =\frac{2a^3}{3\pi},
 \qquad
 S_0(a)=\sum_{n=1}^N b_n^2,
 \qquad D(a)=I_0(a)-S_0(a).
\]
Direct summation of the continuous-minus-discrete product deficit gives
\begin{equation}
 \frac{D(a)}{\pi^2}=\frac{N^2}{2}
 +N\left(t^2+\frac16\right)+\frac{2t^3}{3}.
 \label{cm:eq-optical-product-deficit}
\end{equation}
For \(C_D=1382/3125\), the difference
\(D/\pi^2-C_D(N+t)^2\) increases in \(N\), and at \(N=1\) is minimized
at \(t=C_D\), where it equals
\(1822532/91552734375>0\).  Thus \(D>C_Da^2\).

For \(x>0\),
\[
\left\lceil\sqrt{x^2+\frac14}-\frac12\right\rceil^2
 <x^2+x+1.
\]
Indeed the ceiling is strictly smaller than
\(\sqrt{x^2+1/4}+1/2\), whose square is
\(x^2+1/2+\sqrt{x^2+1/4}<x^2+x+1\).  Applying this with
\(x=b_n/\varepsilon\), then using the monotone left-sum estimate and the
quarter-circle integral,
\[
 \sum_{n=1}^N b_n
 <\int_0^{a/\pi}\sqrt{a^2-\pi^2x^2}\,dx=\frac{a^2}{4},
 \qquad N<\frac a\pi,
\]
gives
\[
 \varepsilon^2\mathcal P_\varepsilon
 <S_0+\varepsilon a^2/4+\varepsilon^2a/\pi.
\]
On \(a\leq c/\varepsilon\), using \(q=106/333>1/\pi\) and
\(\varepsilon\leq11/50\), the exact reserve is
\begin{equation}
 C_D-\left(\frac{2cq}{3}+\frac{11}{200}
 +(11/50)^2q^2\right)
 =\frac{39569}{2772225000}>0.
 \label{cm:eq-optical-low-reserve}
\end{equation}
Indeed, \(a\leq c/\varepsilon\), \(\varepsilon\leq11/50\),
\(a>\pi\), and \(q>1/\pi\) give
\[
 \varepsilon I_0+\frac{\varepsilon a^2}{4}
 +\frac{\varepsilon^2a}{\pi}
 <a^2\left(\frac{2cq}{3}+\frac{11}{200}
 +(11/50)^2q^2\right)<C_Da^2<D.
\]
Consequently
\[
 \varepsilon^2\mathcal P_\varepsilon
 <I_0-D+\frac{\varepsilon a^2}{4}
 +\frac{\varepsilon^2a}{\pi}
 <(1-\varepsilon)I_0<m_\varepsilon I_0
 =\varepsilon^2W,
\]
where the last identity is \eqref{thin:eq-weyl-normalization}.
This proves the low screen, including all radial and angular equality walls.

For the complementary-action screen, use the inverse-action construction
proved in Lemma~\ref{thin:lem-action-geometry}: let the decreasing shell
action have inverse \(R_{\mathcal A}\), put \(F=R_{\mathcal A}^2\), and
\(g=-F'\).  If \(\tau\) is the actual curvature interface and
\(T=a/\pi\), that lemma gives
\[
 g\text{ decreasing on }(0,\tau),\quad
 g\text{ increasing on }(\tau,T),
\]
\[
 F(0)=a^2,\quad F(\tau)=\rho^2a^2,\quad F(T)=0,\quad
 g(T)=2\pi\rho a.
\]
For \(C(t)=\#\{n:n-1/4<t\}\), set \(w=t-C(t)\) and
\(\Psi(t)=\int_0^tw(s)\,ds-t/4\).  The complete sawtooth and Stieltjes
calculation in Lemma~\ref{thin:lem-radial-deficit}, including the ungridded
curvature interface \(\tau\), gives
\begin{equation}
 D_{\rm rad}:=\int_0^TF(t)\,dt
 -\sum_{n-1/4<T}F(n-1/4)
 \geq\frac{\rho^2a^2}{4}-\frac{\pi\rho a}{4}.
 \label{cm:eq-optical-radial-deficit}
\end{equation}
The exact half-integer angular count from
Lemma~\ref{thin:lem-layer-cake} satisfies
\(\varepsilon^2M_\varepsilon(x)^2<(x+\varepsilon/2)^2\), even on an
angular wall.  Since fewer than \(a/\pi+1/4\) radial layers occur, the
angular error is strictly below
\begin{equation}
 E=\left(\frac a\pi+\frac14\right)
 \left(\varepsilon a+\frac{\varepsilon^2}{4}\right).
 \label{cm:eq-optical-angular-error}
\end{equation}
On \(a\geq c/\varepsilon\), one has \(1/a\leq\varepsilon/c\).  Put
\[
 L_{\rm opt}(\varepsilon)
 :=\frac{(1-\varepsilon)^2}{4}
 -\frac{(22/7)(1-\varepsilon)\varepsilon}{4c},
\]
\[
 A_{\rm opt}(\varepsilon)
 :=q\varepsilon+\frac{\varepsilon^2}{4c}
 +\frac{q\varepsilon^3}{4c}
 +\frac{\varepsilon^4}{16c^2}.
\]
Dividing \eqref{cm:eq-optical-radial-deficit} and
\eqref{cm:eq-optical-angular-error} by \(a^2\), then using
\(\pi<22/7\) and \(q>1/\pi\), gives
\[
 \frac{D_{\rm rad}}{a^2}\geq L_{\rm opt}(\varepsilon),
 \qquad
 \frac{E}{a^2}<A_{\rm opt}(\varepsilon).
\]
For \(0<\varepsilon\leq11/50<1/2\),
\[
 L_{\rm opt}'(\varepsilon)
 =-\frac{1-\varepsilon}{2}
 -\frac{22/7}{4c}(1-2\varepsilon)<0,
 \qquad
 A_{\rm opt}'(\varepsilon)>0.
\]
At \(\varepsilon=11/50\), the exact endpoint difference is
\begin{equation}
 L_{\rm opt}(11/50)-A_{\rm opt}(11/50)
 =\frac{14817541}{472867032960000}>0.
 \label{cm:eq-optical-high-reserve}
\end{equation}
Thus the strict layer-cake sum is below the action integral, and the phase
bridge proves the high screen.  The common face lies in the nonzero radial
range because
\[
 \frac c\varepsilon\geq\frac{c}{11/50}
 =\frac{2252}{275}>\frac{22}{7}>\pi,
 \qquad
 \frac{2252}{275}-\frac{22}{7}=\frac{9714}{1925}>0.
\]
Both screens include \(a=c/\varepsilon\), so their union has neither a gap
nor an unowned equality face.  At \(K=0\) both sides vanish.
\end{proof}

% The former residual subtraction and final-staircase closure are archived
% below but no longer enter the revised proof.
\iffalse
Lemma~\ref{cm:lem-optical} owns the outer strip
\(39/50\leq\rho<1\).  Within \(\mathcal D_{19}\), combining
Lemmas~\ref{cm:lem-lower-closure} and \ref{cm:lem-k11} leaves exactly
\begin{equation}
 \boxed{
 \mathcal D_{20}=
 \left\{(\rho,K):\rho_c\leq\rho<\frac{39}{50},\quad
 k_{11}(\rho)<K<U(\rho)=K_0(\rho)\right\}.}
 \label{cm:eq-D20}
\end{equation}
Indeed \(k_{11}<U\): on \(\rho\geq\rho_c\),
\(a_\rho\geq1\) and \(\eta_\rho<1/8\), so \(K_0>64\), whereas
\(k_{11}<64\) for \(\rho<39/50\).  Moreover
\(\omega_0<1/9<7/51<\rho_c\), so \(H_0\) is not eligible on the
surviving ratio range and \(U=K_0\).  The face \(K=k_{11}\) belongs to the closed
staircase, \(\rho=39/50\) belongs to the optical theorem, and \(K=U\)
belongs to the fixed-ratio high-frequency owner.

\subsection{All-frequency analytic closure of the exact residual
\texorpdfstring{\(\mathcal D_{20}\)}{D20}}

The final residual is closed without a frequency split.  The retained-
remainder identity and angular-block estimate are global, while the scalar
gap is strictly increasing in \(K\) above its positive base curve.  The
detailed calculations are printed in Parts~VII--IX of the accompanying
\emph{Purely Analytic Supplement}; the logical assembly is included here.

\begin{lemma}[All-frequency analytic middle theorem]
\label{cm:lem-compact-certificate}
For
\[
 \rho_c\leq\rho\leq\frac{39}{50},
 \qquad
 K\geq k_{11}(\rho),
\]
one has \(N_D(A_{\rho,1},K^2)<W(\rho,K)\).
\end{lemma}

\begin{proof}
Let \(P_{\rm coarse}\) be the strict lattice proxy in
\eqref{cm:eq-coarse}, and let
\[
 N=\#\{n\geq1:n-1/4<(1-\rho)K/\pi\}.
\]
Parts~VII--VIII of the \emph{Purely Analytic Supplement} combine the exact
retained-radial identity with one disjoint left block per active layer.  The
result is the single direct estimate
\begin{equation}
 W-P_{\rm coarse}
 >\mathcal H-\frac{(1-\rho^2)K^2}{6}+\frac N2.
 \label{cm:eq-analytic-direct-gap}
\end{equation}
Here \(\mathcal H\) is the explicit smooth action expression in the
``Direct retained-remainder bound'' of Part~VII.
The strict angular convention is used in the elementary inequality
\(M(r)^2<(r+1/2)^2\), so \eqref{cm:eq-analytic-direct-gap} already includes
every radial--angular common wall.  No separate common-jump lemma is used.

On the displayed closed domain,
\[
 \frac{(1-\rho)k_{11}(\rho)}{\pi}
 =\sqrt{1+\frac{132(1-\rho)^2}{\pi^2}}>1,
\]
so \(N\geq1\).
The remaining scalar theorem in Part~IX is proved on the base curve
\(K=k_{11}(\rho)\) by rational secant majorants and positive Bernstein
coefficients.  Its derivative is strictly positive for every \(K>12\).
Since \(k_{11}(\rho)>241/20>12\), upward integration gives, for every
\(K\geq k_{11}(\rho)\),
\begin{equation}
 \mathcal H-\frac{(1-\rho^2)K^2}{6}+\frac12>0.
 \label{cm:eq-analytic-scalar-bound}
\end{equation}
Since \(N/2\geq1/2\), equations
\eqref{cm:eq-analytic-direct-gap}--\eqref{cm:eq-analytic-scalar-bound} give
\(P_{\rm coarse}<W\).  The phase reduction
\eqref{cm:eq-coarse} proves the claim.  Every inequality used in
\eqref{cm:eq-analytic-direct-gap}--\eqref{cm:eq-analytic-scalar-bound}
is derived in Parts~VII--IX by integration, exact rational comparison,
or positive-side squaring.  No estimate has an upper-frequency premise.
\end{proof}


\begin{proposition}[Exact analytic closure of \(\mathcal D_{20}\)]
\label{cm:prop-exact-closure}
Every point of the compact-middle residual is covered by
Lemma~\ref{cm:lem-compact-certificate}.
\end{proposition}

\begin{proof}
The exact residual description \eqref{cm:eq-D20} gives
\[
 \mathcal D_{20}
 \subset
 \left\{(\rho,K):
 \rho_c\leq\rho\leq\frac{39}{50},\quad
 K\geq k_{11}(\rho)\right\}.
\]
The containing set is precisely the closed domain of
Lemma~\ref{cm:lem-compact-certificate}.  Hence no residual point remains.
The inherited faces \(\rho=39/50\), \(K=k_{11}(\rho)\), and
\(K=U(\rho)\) remain outside \(\mathcal D_{20}\), while
\(\rho=\rho_c\) is included subject to its two strict frequency
inequalities.  The exact owner subtraction is proved row by row in
Parts~I--VI of the \emph{Purely Analytic Supplement}.
\end{proof}

\begin{proposition}[Compact-middle P\'olya theorem]
\label{prop:compact-middle}
For every
\[
 \rho_*<\rho<\frac{39}{50},\qquad K\geq0,
\]
the Dirichlet counting function of the concentric spherical shell satisfies
\[
 N_D(A_{\rho,1},K^2)\leq W(\rho,K).
\]
\end{proposition}

\begin{proof}
Proposition~\ref{cm:prop-primitive-owners} proves the inequality outside
the exact residual \(\mathcal D_{16}\) in \eqref{cm:eq-D16}.  Within that
residual, Lemmas~\ref{cm:lem-k6}, \ref{cm:lem-lower-d},
\ref{cm:lem-lower-closure}, and \ref{cm:lem-k11} remove precisely the regions
subtracted in the definition of \(\mathcal D_{20}\).  Hence any point not
yet owned belongs to
\(\mathcal D_{20}\) in \eqref{cm:eq-D20}.  The analytic compact theorem proves
the inequality on all of \(\mathcal D_{20}\) by
Proposition~\ref{cm:prop-exact-closure}.  Thus no residual point remains.
The cases on every threshold face have
the explicit owners specified above; in particular, equality at a
frequency wall is not lost in taking complements.  Finally, \(K=0\) is
immediate from \(N_D(A_{\rho,1},0)=W(\rho,0)=0\).  This proves the assertion for
the whole stated parameter range.
\end{proof}

\subsection{Analytic dependencies and optional cross-checks}
\label{cm:sec-certificate-protocol}

The logical dependency graph is summarized here before any historical or
reproducibility remarks.  The final column uses only the two statuses shown.
\begin{center}
\scriptsize
\renewcommand{\arraystretch}{1.14}
\setlength{\tabcolsep}{4pt}
\begin{tabular}{@{}p{2.7cm}p{3.5cm}p{3.5cm}p{2.0cm}p{2.2cm}@{}}
\toprule
result&internal premises&external result&support location&logical status\\
\midrule
separated spectrum and exact phase count&
channel forms, determinant, phase monotonicity&
standard Sturm--Liouville and Bessel facts&
main \S2&
used in final proof\\
phase-to-lattice comparison&
exact phase count and zero extension&
FLPS phase enclosure; annuli Theorem~5.1 and Lemma~5.2&
main \S2&
used in final proof\\
weighted scaffold and primitive owners&
multiplicity identity, action integral, aggregate clause&
balls Theorem~5.1; annuli Proposition~3.1 and Theorems~1.4, 3.4&
main \S2--\S4&
used in final proof\\
zero and index registry&
recurrences, rational signs, fixed-channel min--max&
DLMF (10.21.2), (10.47.3); Lorch abstract inequalities~(i), (ii)&
Parts I--III&
used in final proof\\
lower and high staircases&
zero registry, block caps, closed post-event domination&
none beyond the preceding row&
Parts I--IV&
used in final proof\\
optical theorem&
product deficit, radial deficit, angular layer count&
none&
Part V&
used in final proof\\
all-frequency middle theorem&
retained-remainder identity, paired angular bound, scalar positivity&
none&
Parts VII--IX&
used in final proof\\
regional subtraction and closure&
primitive owners and all preceding regional lemmas&
none&
Parts VI and X&
used in final proof\\
retired seam, \(K=200\), aggregate, and executable ledgers&
independent replay contracts only&
Arb only in the retired interval replay&
Part IV and archive/legacy&
optional cross-check\\
\bottomrule
\end{tabular}
\end{center}

No floating-point computation, interval enclosure, or executable truth value
is a premise of the proof.  The finite portion consists of exact identities,
rational cross-multiplications, positive-side squarings, and explicit sign
decompositions.  Parts~I--VI of the \emph{Purely Analytic Supplement} print
the exact zero, threshold, staircase, optical, and owner-subtraction tables,
together with their normalized row census.  Rows attached to the retired
seam and intermediate staircase owners are retained there only as independent
cross-checks; they are not premises of the simplified owner graph.

Parts~VII--IX prove the all-frequency analytic middle theorem by the exact
retained-remainder identity, the paired angular-block inequality, and an
upward-monotone scalar gap.  Together with Parts~I--VI, these arguments
replace the former compact and aggregate certificates and the executable
state-classifier in the logical dependency graph.  The former aggregate
curvature argument is retained separately only as an optional independent
cross-check; it is not a premise of the theorem.

The earlier exact-arithmetic and interval programs are retained separately
as optional regression tools.  Agreement with them is useful evidence
against transcription errors, but their outputs, precision settings, and
hashes are not mathematical assumptions and are not used in any implication
above.

\fi

\section{Proof of the spectral theorem}

We now assemble the three disjoint phase-space owners.  Fix
\(0<\rho<1\) and put
\[
 K_0(\rho):=\frac{\pi}{1-\rho}.
\]
If \(0\leq K\leq K_0(\rho)\), the strict product majorant in
Lemma~\ref{thin:lem-product-majorant} gives
\begin{equation}
 N_D(A_{\rho,1},K^2)=0.
 \label{eq:global-no-mode}
\end{equation}
This includes \(K=K_0(\rho)\): the first one-dimensional Dirichlet
frequency lies on the spectral wall and is excluded.

It remains to take \(K>K_0(\rho)\).  If
\(0<\rho<39/50\), apply
Theorem~\ref{rs:thm:global-proxy}.  If
\(39/50\leq\rho<1\), apply Lemma~\ref{cm:lem-optical}.  These cases,
together with \eqref{eq:global-no-mode}, are disjoint and exhaustive.
Therefore
\begin{equation}
  N_D(A_{\rho,1},K^2)
  \le\frac{2}{9\pi}(1-\rho^3)K^3
  \qquad(K\ge0).
  \label{eq:unit-shell-theorem}
\end{equation}

Now
\[
  \omega_3=\frac{4\pi}{3},
  \qquad
  L_3=\frac{\omega_3}{(2\pi)^3}=\frac1{6\pi^2},
  \qquad
  |A_{\rho,1}|=\frac{4\pi}{3}(1-\rho^3),
\]
so
\begin{equation}
  L_3|A_{\rho,1}|=\frac{2}{9\pi}(1-\rho^3).
  \label{eq:weyl-normalization}
\end{equation}
Taking \(K=\sqrt\Lambda\) in \eqref{eq:unit-shell-theorem} proves the
unit-shell form of Theorem~\ref{thm:spectral} for every \(\Lambda\ge0\).

For arbitrary \(0<r<R\), put \(\rho=r/R\).  The unitary dilation
\[
  (D_Ru)(x)=R^{-3/2}u(x/R)
\]
maps \(L^2(A_{\rho,1})\) onto \(L^2(A_{r,R})\), and its Dirichlet form
scales by \(R^{-2}\).  Hence
\[
  \lambda_j^D(A_{r,R})=R^{-2}\lambda_j^D(A_{\rho,1}),
\]
and the strict threshold transforms as
\begin{equation}
  N_D(A_{r,R},\Lambda)
  =N_D(A_{\rho,1},R^2\Lambda).
  \label{eq:dilation-count}
\end{equation}
Applying the unit-shell result at \(R^2\Lambda\),
\[
\begin{aligned}
  N_D(A_{r,R},\Lambda)
  &\le\frac{2}{9\pi}(1-\rho^3)(R^2\Lambda)^{3/2}\\
  &=\frac{2}{9\pi}(R^3-r^3)\Lambda^{3/2}\\
  &=L_3|A_{r,R}|\Lambda^{3/2}.
\end{aligned}
\]
This proves Theorem~\ref{thm:spectral} for every \(R>0\), including
\(R<1\).

\section{Proof that spherical shells do not tile}

We prove Theorem~\ref{thm:nontiling} independently of the spectral argument.
Set \(\alpha=r/R\in(0,1)\) and scale by \(1/R\).  It suffices to consider
\[
  T=A_{\alpha,1}=\{x:\alpha<|x|<1\}.
\]
Every rigid motion has the form \(x\mapsto Qx+a\), with \(Q\) orthogonal.
Because \(T\) is radial, \(Q(T)=T\), so every congruent copy is a translate
\[
  T_a=a+T=\{x:\alpha<|x-a|<1\}.
\]

Assume for contradiction that \(\{T_a:a\in\mathfrak C\}\) has
pairwise-disjoint interiors and covers almost everywhere.  Duplicate centers
are impossible because they give the same nonempty open tile.

Fix a unit vector \(e\) and put
\[
  c=\frac{1+\alpha}{2}e,
  \qquad
  \delta=\frac{1-\alpha}{2}.
\]
Write \(B(x,s)=\{y\in\R^3:|y-x|<s\}\) for the open Euclidean ball.
Then \(B(c,\delta)\subset T\).  Therefore \(T_a\) contains
\(B(a+c,\delta)\), and disjointness implies
\begin{equation}
  |a-b|\ge1-\alpha
  \qquad(a\ne b).
  \label{eq:center-separation}
\end{equation}
Thus the centers are locally finite: a bounded set cannot contain infinitely
many points with fixed positive separation.  They are also countable, since
\[
  \mathfrak C
  =\bigcup_{m\ge1}(\mathfrak C\cap\overline B(0,m))
\]
is a countable union of finite sets.

Each tile boundary is a union of two spheres and is three-dimensional null.
The countable union
\[
  Z=\bigcup_{a\in\mathfrak C}\partial T_a
\]
is therefore null.  Exact open-shell coverage is already a special case of
almost-everywhere coverage.  If closed shells cover exactly or almost
everywhere, removing \(Z\) shows that their open interiors still cover almost
everywhere.  Hence it is enough to assume
\begin{equation}
  \left|\R^3\setminus\bigcup_{a\in\mathfrak C}T_a\right|=0.
  \label{eq:ae-coverage}
\end{equation}

Choose one tile \(T_{a_0}\) and its outer sphere
\[
  \Sigma=\{x:|x-a_0|=1\}.
\]
Fix \(p\in\Sigma\) and put \(u=p-a_0\).  For \(n\ge1\), let
\[
  U_n=B\left(p+\frac2n u,\frac1n\right).
\]
Every \(U_n\) lies outside \(T_{a_0}\), is open, and has positive volume.
By \eqref{eq:ae-coverage}, choose \(x_n\in U_n\cap T_{a_n}\) with
\(a_n\ne a_0\).  Since \(x_n\to p\) and \(|x_n-a_n|<1\), the centers
\(a_n\) eventually lie in a fixed bounded set.  Local finiteness lets us
pass to a subsequence with one fixed center \(a\ne a_0\).  Thus
\(p\in\overline{T_a}\).

The point \(p\) is not in \(T_a\).  Otherwise some
\(B(p,\varepsilon)\) lies in \(T_a\).  Choose
\[
  0<t<\min\{\varepsilon,1-\alpha\}.
\]
Then \(p-tu\in T_{a_0}\cap T_a\), contradicting disjoint interiors.
Therefore every \(p\in\Sigma\) belongs to the boundary of a neighboring
tile.

If \(p\in\Sigma\cap\partial T_a\), then
\(|p-a|\in\{\alpha,1\}\), hence \(|a-a_0|\le2\).  Only finitely many
centers satisfy this bound.  Consequently \(\Sigma\) is covered by finitely
many neighboring inner or outer boundary spheres.

Translate so that \(a_0=0\).  The intersection of \(\Sigma\) with a
neighboring sphere \(\{x:|x-a|=s\}\), \(s\in\{\alpha,1\}\), satisfies
\[
  |x|^2=1,
  \qquad |x-a|^2=s^2,
\]
and hence
\begin{equation}
  2a\cdot x=|a|^2+1-s^2.
  \label{eq:sphere-plane}
\end{equation}
Because \(a\ne0\), \eqref{eq:sphere-plane} is a genuine plane.  Its
intersection with \(\Sigma\) is empty, a tangent point, or a circle, and
therefore has zero surface measure on \(\Sigma\).  A finite union of such
sets cannot cover \(\Sigma\), whose area is \(4\pi\).

There is no coincident-sphere exception.  A neighboring outer sphere equals
\(\Sigma\) only when it has the same center, which is the original tile or a
forbidden duplicate.  A neighboring inner sphere cannot equal \(\Sigma\)
because \(\alpha<1\).  We have reached a contradiction.  Scaling back proves
Theorem~\ref{thm:nontiling}.

\appendix

\section{External inputs and verification boundary}
\label{app:external-boundary}

The proof above is self-contained with respect to every shell-specific
spectral reduction, weighted count, endpoint estimate, regional subtraction,
fixed-channel variational comparison, boundary assignment, and finite
exact-rational predicate.  It uses the following conventional external results, each only
in the scope stated where it enters the proof:
\begin{enumerate}
  \item standard spherical-harmonic decomposition, regular
        Sturm--Liouville theory, the contiguous-order zero interlacing
        in DLMF equation~(10.21.2), and the classical Bessel identities
        collected in \cite{DLMF};
  \item the uniform one-variable Bessel-phase enclosure of
        \cite{FLPSphase}, together with the explicitly quoted floor-sum
        inequalities of \cite{FLPSballs,FLPSannuli};
  \item inequalities~(i) and (ii), the two strict lower bounds in Lorch's
        published abstract, for the first positive zero of \(J_\nu\) on
        \(-1<\nu<\infty\), in precisely the half-integer specializations
        derived in the text \cite{Lorch1993}.
\end{enumerate}

All shell-specific finite decisions are exposed in the accompanying
\emph{Purely Analytic Supplement}.  They require only exact integer and
rational arithmetic, elementary integral inequalities, and the classical
real functions already present in the analytic argument.  The archived
interval programs have no place in this verification boundary; they are
optional regression checks only.

\begin{thebibliography}{5}

\bibitem{FLPSballs}
N.~Filonov, M.~Levitin, I.~Polterovich, and D.~A.~Sher,
\newblock P\'olya's conjecture for Euclidean balls,
\newblock \emph{Invent. Math.} \textbf{234} (2023), 129--169.
\newblock \href{https://doi.org/10.1007/s00222-023-01198-1}{doi:10.1007/s00222-023-01198-1}.

\bibitem{FLPSphase}
N.~Filonov, M.~Levitin, I.~Polterovich, and D.~A.~Sher,
\newblock Uniform enclosures for the phase and zeros of Bessel functions and
their derivatives,
\newblock \emph{SIAM J. Math. Anal.} \textbf{56} (2024), no.~6, 7644--7682.
\newblock \href{https://doi.org/10.1137/24M1642032}{doi:10.1137/24M1642032}.

\bibitem{FLPSannuli}
N.~Filonov, M.~Levitin, I.~Polterovich, and D.~A.~Sher,
\newblock P\'olya's conjecture for Dirichlet eigenvalues of annuli,
\newblock \emph{J. Lond. Math. Soc.} \textbf{113} (2026), no.~2, e70425.
\newblock \href{https://doi.org/10.1112/jlms.70425}{doi:10.1112/jlms.70425}.

\bibitem{Lorch1993}
L.~Lorch,
\newblock Some inequalities for the first positive zeros of Bessel
functions,
\newblock \emph{SIAM J. Math. Anal.} \textbf{24} (1993), no.~3, 814--823.
\newblock \href{https://doi.org/10.1137/0524050}{doi:10.1137/0524050}.

\bibitem{DLMF}
NIST Digital Library of Mathematical Functions,
\newblock Chapter 10: Bessel functions, in particular equation (10.9.30),
\S10.21(i), equation (10.21.2), and equation (10.47.3),
\newblock \url{https://dlmf.nist.gov/10}, accessed 16 July 2026.

\end{thebibliography}

\end{document}
